% Автоматаар үүсгэсэн — эх файл: SETUP.md
\chapter{Ажлын орчны бэлтгэл}

Энэ заавраар \textbf{нэг удаа} бэлтгэсэн орчинг 8 лабораторийн ажил бүгд ашиглана. Үүнийг лабораторийн цагаар биш, \textbf{I--II долоо хоногийн бие даалтаар} гүйцэтгэнэ. Лаб 1-д ирэхээс өмнө бүх алхмыг дуусгаж, \textbf{Г хэсгийн хяналтын жагсаалтыг} бүгдийг ✔ болгосон байх ёстой.

\begingroup\scriptsize\begin{longtable}[]{@{}llll@{}}
\toprule
\begin{minipage}[b]{0.22\columnwidth}\raggedright
Хэсэг\strut
\end{minipage} & \begin{minipage}[b]{0.22\columnwidth}\raggedright
Хаана\strut
\end{minipage} & \begin{minipage}[b]{0.22\columnwidth}\raggedright
Юу хийх вэ\strut
\end{minipage} & \begin{minipage}[b]{0.22\columnwidth}\raggedright
Хугацаа\strut
\end{minipage}\tabularnewline
\midrule
\endhead
\begin{minipage}[t]{0.22\columnwidth}\raggedright
\textbf{0}\strut
\end{minipage} & \begin{minipage}[t]{0.22\columnwidth}\raggedright
---\strut
\end{minipage} & \begin{minipage}[t]{0.22\columnwidth}\raggedright
Уншиж, өөрийн утгуудаа тэмдэглэх\strut
\end{minipage} & \begin{minipage}[t]{0.22\columnwidth}\raggedright
10 мин\strut
\end{minipage}\tabularnewline
\begin{minipage}[t]{0.22\columnwidth}\raggedright
\textbf{А}\strut
\end{minipage} & \begin{minipage}[t]{0.22\columnwidth}\raggedright
{[}ЗК{]} Зөөврийн компьютер (үүл)\strut
\end{minipage} & \begin{minipage}[t]{0.22\columnwidth}\raggedright
WSL 2 + Ubuntu, Git, VS Code, Docker Desktop, Raspberry Pi Imager → үүлний стек\strut
\end{minipage} & \begin{minipage}[t]{0.22\columnwidth}\raggedright
1.5--2 цаг\strut
\end{minipage}\tabularnewline
\begin{minipage}[t]{0.22\columnwidth}\raggedright
\textbf{Б}\strut
\end{minipage} & \begin{minipage}[t]{0.22\columnwidth}\raggedright
{[}Pi{]} Raspberry Pi 3B (ирмэг)\strut
\end{minipage} & \begin{minipage}[t]{0.22\columnwidth}\raggedright
microSD бичих, Docker Engine, swap, cgroup, ирмэгийн орчин\strut
\end{minipage} & \begin{minipage}[t]{0.22\columnwidth}\raggedright
1.5--2 цаг\strut
\end{minipage}\tabularnewline
\begin{minipage}[t]{0.22\columnwidth}\raggedright
\textbf{В}\strut
\end{minipage} & \begin{minipage}[t]{0.22\columnwidth}\raggedright
{[}ЗК{]}{[}Pi{]} Хоёулаа\strut
\end{minipage} & \begin{minipage}[t]{0.22\columnwidth}\raggedright
Гүүр холбож, мессеж дамжихыг шалгах\strut
\end{minipage} & \begin{minipage}[t]{0.22\columnwidth}\raggedright
15 мин\strut
\end{minipage}\tabularnewline
\begin{minipage}[t]{0.22\columnwidth}\raggedright
\textbf{Г}\strut
\end{minipage} & \begin{minipage}[t]{0.22\columnwidth}\raggedright
{[}ЗК{]}{[}Pi{]}\strut
\end{minipage} & \begin{minipage}[t]{0.22\columnwidth}\raggedright
Хяналтын жагсаалт\strut
\end{minipage} & \begin{minipage}[t]{0.22\columnwidth}\raggedright
10 мин\strut
\end{minipage}\tabularnewline
\begin{minipage}[t]{0.22\columnwidth}\raggedright
\textbf{Д}\strut
\end{minipage} & \begin{minipage}[t]{0.22\columnwidth}\raggedright
---\strut
\end{minipage} & \begin{minipage}[t]{0.22\columnwidth}\raggedright
Түгээмэл алдаа ба шийдэл\strut
\end{minipage} & \begin{minipage}[t]{0.22\columnwidth}\raggedright
хэрэгтэй үед\strut
\end{minipage}\tabularnewline
\bottomrule
\end{longtable}\endgroup

Нийт \textbf{3--4 цаг}, ихэнх нь татаж авах хугацаа --- \textbf{хурдан интернэттэй газар} хий. Лаб 8-ын K3s server VM (А.8)-ийг Лаб 8-аас өмнө тусад нь бэлтгэнэ (+30 мин).

\begin{quote}
☞ \textbf{Алхмуудыг дарааллаар нь хий.} Алхам бүрийн төгсгөлд ✓ \textbf{Шалгах} гэсэн хэсэг бий. Тэнд бичсэн үр дүн гарахгүй бол \textbf{цааш бүү яв} --- Д хэсгээс шалтгааныг хайж ол. Алдааг дараа нь засах нь хамаагүй хэцүү.
\end{quote}

\hypertarget{ux44dux445ux43bux44dux445ux44dux44dux441-ux4e9ux43cux43dux4e9-ux443ux43dux448ux438ux43dux430-ux443ux443}{%
\section{0. Эхлэхээс өмнө уншина уу}\label{ux44dux445ux43bux44dux445ux44dux44dux441-ux4e9ux43cux43dux4e9-ux443ux43dux448ux438ux43dux430-ux443ux443}}

\hypertarget{ux43aux43eux43cux430ux43dux434ux44bux433-ux445ux430ux430ux43dux430-ux431ux438ux447ux438ux445ux438ux439ux433-ux430ux43dux445ux430ux430ux440}{%
\subsection{0.1 Командыг ХААНА бичихийг анхаар}\label{ux43aux43eux43cux430ux43dux434ux44bux433-ux445ux430ux430ux43dux430-ux431ux438ux447ux438ux445ux438ux439ux433-ux430ux43dux445ux430ux430ux440}}

Энэ курст \textbf{гурван өөр терминал} ашиглана. Код блок бүрийн дээр командыг аль терминалд бичихийг заасан:

\begingroup\scriptsize\begin{longtable}[]{@{}llll@{}}
\toprule
\begin{minipage}[b]{0.22\columnwidth}\raggedright
Тэмдэглэгээ\strut
\end{minipage} & \begin{minipage}[b]{0.22\columnwidth}\raggedright
Хаана ажиллана\strut
\end{minipage} & \begin{minipage}[b]{0.22\columnwidth}\raggedright
Хэрхэн нээх вэ\strut
\end{minipage} & \begin{minipage}[b]{0.22\columnwidth}\raggedright
Prompt (мөрийн эхлэл)\strut
\end{minipage}\tabularnewline
\midrule
\endhead
\begin{minipage}[t]{0.22\columnwidth}\raggedright
\textbf{{[}ЗК · PowerShell{]}}\strut
\end{minipage} & \begin{minipage}[t]{0.22\columnwidth}\raggedright
Windows\strut
\end{minipage} & \begin{minipage}[t]{0.22\columnwidth}\raggedright
Start → \texttt{Terminal} (Windows 10-д \texttt{PowerShell})\strut
\end{minipage} & \begin{minipage}[t]{0.22\columnwidth}\raggedright
\texttt{PS\ C:\textbackslash{}Users\textbackslash{}bat\textgreater{}}\strut
\end{minipage}\tabularnewline
\begin{minipage}[t]{0.22\columnwidth}\raggedright
\textbf{{[}ЗК · PowerShell (Admin){]}}\strut
\end{minipage} & \begin{minipage}[t]{0.22\columnwidth}\raggedright
Windows, администраторын эрхтэй\strut
\end{minipage} & \begin{minipage}[t]{0.22\columnwidth}\raggedright
Start → \texttt{Terminal} → \textbf{баруун товч} → \textbf{Run as administrator} → \textbf{Yes}\strut
\end{minipage} & \begin{minipage}[t]{0.22\columnwidth}\raggedright
\texttt{PS\ C:\textbackslash{}WINDOWS\textbackslash{}system32\textgreater{}}\strut
\end{minipage}\tabularnewline
\begin{minipage}[t]{0.22\columnwidth}\raggedright
\textbf{{[}ЗК · Ubuntu{]}}\strut
\end{minipage} & \begin{minipage}[t]{0.22\columnwidth}\raggedright
Зөөврийн компьютер доторх Linux (WSL)\strut
\end{minipage} & \begin{minipage}[t]{0.22\columnwidth}\raggedright
Start → \texttt{Ubuntu} (эсвэл Terminal-ын \texttt{▼} товч → \textbf{Ubuntu})\strut
\end{minipage} & \begin{minipage}[t]{0.22\columnwidth}\raggedright
\texttt{bat@LAPTOP:\textasciitilde{}\$}\strut
\end{minipage}\tabularnewline
\begin{minipage}[t]{0.22\columnwidth}\raggedright
\textbf{{[}Pi{]}}\strut
\end{minipage} & \begin{minipage}[t]{0.22\columnwidth}\raggedright
Raspberry Pi дээр, SSH-ээр\strut
\end{minipage} & \begin{minipage}[t]{0.22\columnwidth}\raggedright
Б.2-т заасан\strut
\end{minipage} & \begin{minipage}[t]{0.22\columnwidth}\raggedright
\texttt{cnc302@pi-team07:\textasciitilde{}\ \$}\strut
\end{minipage}\tabularnewline
\bottomrule
\end{longtable}\endgroup

\begin{quote}
⚠ \textbf{Хамгийн түгээмэл алдаа:} командыг буруу терминалд бичих. Командыг бичихээсээ өмнө \textbf{prompt-ыг хар}.

☞ \textbf{Лабораторийн зааварт} {[}ЗК{]} гэж тэмдэглэсэн бүх командыг \textbf{{[}ЗК · Ubuntu{]}} терминалд бичнэ. Лабын заавар \texttt{make}, \texttt{mosquitto\_sub}, \texttt{python3} зэрэг Linux командыг ашигладаг --- тэдгээр нь Windows PowerShell-д байхгүй, харин Ubuntu (WSL) дотор шууд ажиллана. PowerShell-ийг зөвхөн энэ SETUP-д, Windows-ийн тохиргоонд л хэрэглэнэ.
\end{quote}

\begin{quote}
☞ \textbf{macOS / Linux зөөврийн компьютертэй бол:} WSL-тэй холбоотой алхмуудыг (А.1.3, А.2-ын \texttt{.wslconfig}) алгас. {[}ЗК · PowerShell{]} ба {[}ЗК · Ubuntu{]} гэсэн бүх командыг өөрийн \textbf{Terminal}-д бич. Docker Desktop-ийн санах ойг \textbf{Settings → Resources → Advanced → Memory limit}-ээр тохируулна (А.2).
\end{quote}

\hypertarget{ux43aux43eux43cux430ux43dux434ux44bux433-ux445ux44dux440ux445ux44dux43d-ux445ux443ux443ux43bux430ux445-ux432ux44d}{%
\subsection{0.2 Командыг хэрхэн хуулах вэ}\label{ux43aux43eux43cux430ux43dux434ux44bux433-ux445ux44dux440ux445ux44dux43d-ux445ux443ux443ux43bux430ux445-ux432ux44d}}

\begin{enumerate}
\def\labelenumi{\arabic{enumi}.}
\tightlist
\item
  Код блокийн баруун дээд буланд байгаа \textbf{хуулах (copy)} товчийг дар --- эсвэл хулганаар бүтэн мөрийг сонгож \texttt{Ctrl\ +\ C}.
\item
  Терминал дээр \textbf{баруун товч} дар (эсвэл \texttt{Ctrl\ +\ Shift\ +\ V}) --- буулгана.
\item
  \texttt{Enter} дар.
\item
  Команд \textbf{дуусахыг хүлээ}: prompt дахин гарч ирсний дараа дараагийн командаа бич.
\end{enumerate}

\begin{itemize}
\tightlist
\item
  \texttt{sudo} гэж эхэлсэн команд \textbf{нууц үг} асууна. Нууц үг бичихэд дэлгэцэнд \textbf{юу ч гарахгүй} (одоор ч харагдахгүй) --- энэ хэвийн. Бичээд \texttt{Enter} дар.
\item
  \texttt{\#}-ээс хойшхи хэсэг бол \textbf{тайлбар} --- бичих шаардлагагүй.
\item
  Командыг зогсоох: \texttt{Ctrl\ +\ C}.
\end{itemize}

\hypertarget{ux4e9ux4e9ux440ux438ux439ux43d-ux443ux442ux433ux443ux443ux434ux44bux433-ux442ux44dux43cux434ux44dux433ux43bux44d}{%
\subsection{0.3 Өөрийн утгуудыг тэмдэглэ}\label{ux4e9ux4e9ux440ux438ux439ux43d-ux443ux442ux433ux443ux443ux434ux44bux433-ux442ux44dux43cux434ux44dux433ux43bux44d}}

Заавар дахь \texttt{\textless{}\ldots{}\textgreater{}} хэсгийг өөрийн утгаар солино. \textbf{Хаалтыг (\texttt{\textless{}\ \textgreater{}}) хамт устгана.} Жишээ: \texttt{pi-team\textless{}NN\textgreater{}} → \texttt{pi-team07}.

\begingroup\scriptsize\begin{longtable}[]{@{}llll@{}}
\toprule
\begin{minipage}[b]{0.22\columnwidth}\raggedright
Хувьсагч\strut
\end{minipage} & \begin{minipage}[b]{0.22\columnwidth}\raggedright
Утга\strut
\end{minipage} & \begin{minipage}[b]{0.22\columnwidth}\raggedright
Жишээ\strut
\end{minipage} & \begin{minipage}[b]{0.22\columnwidth}\raggedright
Хаанаас\strut
\end{minipage}\tabularnewline
\midrule
\endhead
\begin{minipage}[t]{0.22\columnwidth}\raggedright
\texttt{\textless{}NN\textgreater{}}\strut
\end{minipage} & \begin{minipage}[t]{0.22\columnwidth}\raggedright
Багийн дугаар, \textbf{хоёр оронтой}\strut
\end{minipage} & \begin{minipage}[t]{0.22\columnwidth}\raggedright
\texttt{07}\strut
\end{minipage} & \begin{minipage}[t]{0.22\columnwidth}\raggedright
Багш\strut
\end{minipage}\tabularnewline
\begin{minipage}[t]{0.22\columnwidth}\raggedright
\texttt{\textless{}REPO\_\allowbreak{}URL\textgreater{}}\strut
\end{minipage} & \begin{minipage}[t]{0.22\columnwidth}\raggedright
Багийн Git сангийн хаяг\strut
\end{minipage} & \begin{minipage}[t]{0.22\columnwidth}\raggedright
\texttt{https:\allowbreak{}/\allowbreak{}/\allowbreak{}github.\allowbreak{}com/\allowbreak{}\ldots{}/\allowbreak{}cnc302-team07.\allowbreak{}git}\strut
\end{minipage} & \begin{minipage}[t]{0.22\columnwidth}\raggedright
Багш\strut
\end{minipage}\tabularnewline
\begin{minipage}[t]{0.22\columnwidth}\raggedright
\texttt{\textless{}LAPTOP\_\allowbreak{}IP\textgreater{}}\strut
\end{minipage} & \begin{minipage}[t]{0.22\columnwidth}\raggedright
Зөөврийн компьютерийн LAN IP\strut
\end{minipage} & \begin{minipage}[t]{0.22\columnwidth}\raggedright
\texttt{192.168.1.100}\strut
\end{minipage} & \begin{minipage}[t]{0.22\columnwidth}\raggedright
А.5\strut
\end{minipage}\tabularnewline
\begin{minipage}[t]{0.22\columnwidth}\raggedright
\texttt{\textless{}PI\_\allowbreak{}IP\textgreater{}}\strut
\end{minipage} & \begin{minipage}[t]{0.22\columnwidth}\raggedright
Pi-гийн IP хаяг\strut
\end{minipage} & \begin{minipage}[t]{0.22\columnwidth}\raggedright
\texttt{192.168.1.57}\strut
\end{minipage} & \begin{minipage}[t]{0.22\columnwidth}\raggedright
Б.2\strut
\end{minipage}\tabularnewline
\bottomrule
\end{longtable}\endgroup

Pi-гийн hostname нь \texttt{pi-team\textless{}NN\textgreater{}}, хэрэглэгч нь \textbf{\texttt{cnc302}}, ирмэгийн \texttt{DEVICE\_ID} нь \texttt{pi3b-team\textless{}NN\textgreater{}} байна. Эдгээр нэрийг \textbf{өөрчлөхгүй} --- лабын код ба \texttt{edge/\allowbreak{}agent/\allowbreak{}cnc302-edge-agent.\allowbreak{}service} яг эдгээр нэрийг хүлээнэ.

\hypertarget{ux442ux43eux43dux43eux433-ux442ux4e9ux445ux4e9ux4e9ux440ux4e9ux43cux436ux438ux439ux43d-ux448ux430ux430ux440ux434ux43bux430ux433ux430}{%
\subsection{0.4 Тоног төхөөрөмжийн шаардлага}\label{ux442ux43eux43dux43eux433-ux442ux4e9ux445ux4e9ux4e9ux440ux4e9ux43cux436ux438ux439ux43d-ux448ux430ux430ux440ux434ux43bux430ux433ux430}}

\begingroup\scriptsize\begin{longtable}[]{@{}ll@{}}
\toprule
\begin{minipage}[b]{0.47\columnwidth}\raggedright
Зүйл\strut
\end{minipage} & \begin{minipage}[b]{0.47\columnwidth}\raggedright
Шаардлага\strut
\end{minipage}\tabularnewline
\midrule
\endhead
\begin{minipage}[t]{0.47\columnwidth}\raggedright
Зөөврийн компьютер\strut
\end{minipage} & \begin{minipage}[t]{0.47\columnwidth}\raggedright
Windows 10 22H2 (build 19045) эсвэл Windows 11 23H2 (build 22631)+, 64-bit; \textbf{RAM ≥ 8 GB} (16 GB тав тухтай --- Лаб 8); дискэнд \textbf{≥ 40 GB} сул зай\strut
\end{minipage}\tabularnewline
\begin{minipage}[t]{0.47\columnwidth}\raggedright
Raspberry Pi 3B\strut
\end{minipage} & \begin{minipage}[t]{0.47\columnwidth}\raggedright
+ \textbf{5 V / 2.5 A} тэжээлийн адаптер (утасны цэнэглэгч хангалтгүй) + \textbf{microSD 16 GB+} (A1/A2) + microSD уншигч\strut
\end{minipage}\tabularnewline
\begin{minipage}[t]{0.47\columnwidth}\raggedright
Сүлжээ\strut
\end{minipage} & \begin{minipage}[t]{0.47\columnwidth}\raggedright
Зөөврийн компьютер ба Pi \textbf{нэг сүлжээнд}. Pi-г \textbf{кабелиар} холбох нь хамгийн найдвартай (А.5)\strut
\end{minipage}\tabularnewline
\bottomrule
\end{longtable}\endgroup

\hypertarget{ux430.-ux437ux43a-ux437ux4e9ux4e9ux432ux440ux438ux439ux43d-ux43aux43eux43cux43fux44cux44eux442ux435ux440-ux434ux44dux44dux440}{%
\section{А. {[}ЗК{]} Зөөврийн компьютер дээр}\label{ux430.-ux437ux43a-ux437ux4e9ux4e9ux432ux440ux438ux439ux43d-ux43aux43eux43cux43fux44cux44eux442ux435ux440-ux434ux44dux44dux440}}

\hypertarget{ux430.1-ux43fux440ux43eux433ux440ux430ux43cux443ux443ux434-ux441ux443ux443ux43bux433ux430ux445}{%
\subsection{А.1 Програмууд суулгах}\label{ux430.1-ux43fux440ux43eux433ux440ux430ux43cux443ux443ux434-ux441ux443ux443ux43bux433ux430ux445}}

Суулгах дараалал \textbf{чухал}: WSL → Git → VS Code → Docker Desktop → Raspberry Pi Imager. Docker Desktop нь WSL 2 дээр ажилладаг тул WSL-ийг \textbf{эхэлж} суулгана.

\begingroup\scriptsize\begin{longtable}[]{@{}lll@{}}
\toprule
\begin{minipage}[b]{0.30\columnwidth}\raggedright
\#\strut
\end{minipage} & \begin{minipage}[b]{0.30\columnwidth}\raggedright
Програм\strut
\end{minipage} & \begin{minipage}[b]{0.30\columnwidth}\raggedright
Юунд хэрэглэх вэ\strut
\end{minipage}\tabularnewline
\midrule
\endhead
\begin{minipage}[t]{0.30\columnwidth}\raggedright
А.1.3\strut
\end{minipage} & \begin{minipage}[t]{0.30\columnwidth}\raggedright
\textbf{WSL 2 + Ubuntu}\strut
\end{minipage} & \begin{minipage}[t]{0.30\columnwidth}\raggedright
Windows доторх Linux. Лабын бүх {[}ЗК{]} команд энд ажиллана\strut
\end{minipage}\tabularnewline
\begin{minipage}[t]{0.30\columnwidth}\raggedright
А.1.4\strut
\end{minipage} & \begin{minipage}[t]{0.30\columnwidth}\raggedright
\textbf{Git for Windows}\strut
\end{minipage} & \begin{minipage}[t]{0.30\columnwidth}\raggedright
Git-ийн нэвтрэлт (Credential Manager), VS Code-ийн Git\strut
\end{minipage}\tabularnewline
\begin{minipage}[t]{0.30\columnwidth}\raggedright
А.1.5\strut
\end{minipage} & \begin{minipage}[t]{0.30\columnwidth}\raggedright
\textbf{VS Code}\strut
\end{minipage} & \begin{minipage}[t]{0.30\columnwidth}\raggedright
Код засах, Ubuntu (WSL) ба Pi (SSH) руу холбогдох\strut
\end{minipage}\tabularnewline
\begin{minipage}[t]{0.30\columnwidth}\raggedright
А.1.6\strut
\end{minipage} & \begin{minipage}[t]{0.30\columnwidth}\raggedright
\textbf{Docker Desktop}\strut
\end{minipage} & \begin{minipage}[t]{0.30\columnwidth}\raggedright
Үүлний стек (EMQX, InfluxDB, Grafana, \ldots)\strut
\end{minipage}\tabularnewline
\begin{minipage}[t]{0.30\columnwidth}\raggedright
А.1.7\strut
\end{minipage} & \begin{minipage}[t]{0.30\columnwidth}\raggedright
\textbf{Raspberry Pi Imager}\strut
\end{minipage} & \begin{minipage}[t]{0.30\columnwidth}\raggedright
microSD карт бичих\strut
\end{minipage}\tabularnewline
\begin{minipage}[t]{0.30\columnwidth}\raggedright
А.1.9\strut
\end{minipage} & \begin{minipage}[t]{0.30\columnwidth}\raggedright
Ubuntu доторх хэрэгслүүд\strut
\end{minipage} & \begin{minipage}[t]{0.30\columnwidth}\raggedright
\texttt{make}, \texttt{python3-venv}, \texttt{mosquitto-clients}, \texttt{git}, \texttt{curl}\strut
\end{minipage}\tabularnewline
\bottomrule
\end{longtable}\endgroup

\hypertarget{ux430.1.1-windows-ux438ux439ux43d-ux445ux443ux432ux438ux43bux431ux430ux440-ux431ux430-virtualization-ux448ux430ux43bux433ux430ux445}{%
\subsubsection{А.1.1 Windows-ийн хувилбар ба virtualization шалгах}\label{ux430.1.1-windows-ux438ux439ux43d-ux445ux443ux432ux438ux43bux431ux430ux440-ux431ux430-virtualization-ux448ux430ux43bux433ux430ux445}}

\textbf{Алхам 1 --- Windows-ийн хувилбар.} \texttt{Win\ +\ R} → \texttt{winver} гэж бичээд \texttt{Enter}.

✓ \textbf{Шалгах:} Нээгдсэн цонхонд \textbf{Version 22H2} (Windows 10) эсвэл \textbf{Version 23H2 / 24H2 / 25H2} (Windows 11) байна. Хуучин бол \textbf{Settings → Windows Update → Check for updates} → бүх шинэчлэлийг суулгаж, restart хий, дахин шалга.

\begin{quote}
☞ Windows \textbf{Home} хувилбар ч болно. Docker-ийн баримтаар Home нь зөвхөн \emph{Linux контейнер} ажиллуулна --- бидэнд яг тэр л хэрэгтэй.
\end{quote}

\textbf{Алхам 2 --- Virtualization.}

\begin{enumerate}
\def\labelenumi{\arabic{enumi}.}
\tightlist
\item
  \texttt{Ctrl\ +\ Shift\ +\ Esc} → \textbf{Task Manager} нээгдэнэ.
\item
  Зүүн талаас \textbf{Performance} → \textbf{CPU} сонго.
\item
  Баруун доод хэсгээс \textbf{Virtualization}-ийг хай.
\end{enumerate}

✓ \textbf{Шалгах:} \texttt{Virtualization:\allowbreak{}\ Enabled}.

✗ \textbf{Disabled} бол компьютерээ restart хийж, асах үед BIOS руу ор (DELL: \texttt{F2}; бусад: \texttt{F2}, \texttt{Del}, \texttt{F10} эсвэл \texttt{Esc}). \textbf{Intel Virtualization Technology (VT-x)} эсвэл AMD дээр \textbf{SVM Mode}-ийг хайж \textbf{Enabled} болго → \textbf{Save \& Exit}. Олдохгүй бол багшид хандана.

\hypertarget{ux430.1.2-winget-ux431ux430-windows-terminal-ux448ux430ux43bux433ux430ux445}{%
\subsubsection{\texorpdfstring{А.1.2 \texttt{winget} ба Windows Terminal шалгах}{А.1.2 winget ба Windows Terminal шалгах}}\label{ux430.1.2-winget-ux431ux430-windows-terminal-ux448ux430ux43bux433ux430ux445}}

\texttt{winget} бол Windows-ийн багц менежер: програмыг \textbf{албан ёсны эх сурвалжаас} нэг командаар татаж суулгана.

\textbf{{[}ЗК · PowerShell{]}}

\begin{Shaded}
\begin{Highlighting}[]
\NormalTok{winget {-}{-}version}
\end{Highlighting}
\end{Shaded}

✓ \textbf{Шалгах:} \texttt{v1.x.x} гэж хэвлэнэ.

✗ \texttt{winget\ is\ not\ recognized} гэвэл: \textbf{Microsoft Store} → \texttt{App\ Installer} гэж хайж → \textbf{Update} (эсвэл \textbf{Get}) → PowerShell-ээ хаагаад дахин нээ.

Windows 10 дээр \textbf{Windows Terminal} байхгүй бол суулга (Windows 11-д суулгаастай):

\textbf{{[}ЗК · PowerShell{]}}

\begin{Shaded}
\begin{Highlighting}[]
\NormalTok{winget install {-}{-}id Microsoft.}\FunctionTok{WindowsTerminal}\NormalTok{ {-}e}
\end{Highlighting}
\end{Shaded}

\begin{quote}
☞ \texttt{winget}-ийг \textbf{анх удаа} ажиллуулахад \texttt{Do\ you\ agree\ to\ all\ the\ source\ agreements\ terms?\allowbreak{}\ {[}Y{]}\ Yes\ {[}N{]}\ No} гэж асууна → \texttt{Y} дараад \texttt{Enter}. Суулгах явцад Windows \textbf{"Do you want to allow this app to make changes?"} гэж асуувал \textbf{Yes}.
\end{quote}

\hypertarget{ux430.1.3-wsl-2-ux431ux430-ubuntu-ux441ux443ux443ux43bux433ux430ux445}{%
\subsubsection{А.1.3 WSL 2 ба Ubuntu суулгах}\label{ux430.1.3-wsl-2-ux431ux430-ubuntu-ux441ux443ux443ux43bux433ux430ux445}}

WSL (Windows Subsystem for Linux) нь Windows дотор жинхэнэ Linux ажиллуулна. Docker Desktop контейнерүүдээ үүн дотор ажиллуулдаг бөгөөд бид лабын командуудаа \textbf{Ubuntu} дотор бичнэ.

\textbf{Алхам 1 --- суулгах.}

\textbf{{[}ЗК · PowerShell (Admin){]}}

\begin{Shaded}
\begin{Highlighting}[]
\NormalTok{wsl {-}{-}install {-}d Ubuntu}
\end{Highlighting}
\end{Shaded}

Татаж суулгахад 5--10 мин болно. \texttt{The\ requested\ operation\ is\ successful.\allowbreak{}\ Changes\ will\ not\ be\ effective\ until\ the\ system\ is\ rebooted.\allowbreak{}} гэвэл \textbf{компьютерээ restart хий}.

\begin{quote}
✗ Явц \texttt{0.0\%} дээр удаан зогсвол \texttt{Ctrl\ +\ C} дараад: \texttt{wsl\ -\/\allowbreak{}-install\ -\/\allowbreak{}-web-download\ -d\ Ubuntu}

✗ Команд суулгахын оронд тусламжийн текст хэвлэвэл WSL аль хэдийн суусан байна → \texttt{wsl\ -\/-update} хийгээд дараа нь \texttt{wsl\ -\/\allowbreak{}-install\ -d\ Ubuntu}.
\end{quote}

\textbf{Алхам 2 --- Ubuntu-гийн хэрэглэгч үүсгэх.} Restart хийсний дараа \textbf{Ubuntu} цонх автоматаар нээгдэнэ (нээгдэхгүй бол Start → \texttt{Ubuntu}). Хэдэн минут хүлээсний дараа:

\begin{verbatim}
Enter new UNIX username:
\end{verbatim}

→ \textbf{жижиг латин үсгээр}, хоосон зайгүй нэр бич (жишээ: \texttt{bat}) → \texttt{Enter}.

\begin{verbatim}
New password:
Retype new password:
\end{verbatim}

→ нууц үгээ хоёр удаа бич (\textbf{дэлгэцэнд харагдахгүй --- хэвийн}). Энэ бол Ubuntu-гийн \texttt{sudo} нууц үг --- \textbf{тэмдэглэж ав}.

✓ Prompt \texttt{bat@LAPTOP-XXXX:\textasciitilde{}\$} болж өөрчлөгдвөл Ubuntu бэлэн.

\textbf{Алхам 3 --- WSL-ийн хувилбар шалгах.}

\textbf{{[}ЗК · PowerShell{]}}

\begin{Shaded}
\begin{Highlighting}[]
\NormalTok{wsl {-}{-}update}
\NormalTok{wsl {-}{-}version}
\NormalTok{wsl {-}{-}list {-}{-}verbose}
\end{Highlighting}
\end{Shaded}

✓ \textbf{Шалгах:}

\begin{itemize}
\tightlist
\item
  \texttt{wsl\ -\/-version}-ийн эхний мөр: \texttt{WSL\ version:\allowbreak{}\ 2.\allowbreak{}x.\allowbreak{}x} --- \textbf{2.1.5-аас бага биш} (Docker Desktop-ийн шаардлага).
\item
  \texttt{wsl\ -\/\allowbreak{}-list\ -\/\allowbreak{}-verbose}: \texttt{Ubuntu} мөрийн \textbf{VERSION} баганад \texttt{2}, нэрийн өмнө \texttt{*} (анхдагч distribution).
\end{itemize}

✗ VERSION \texttt{1} бол: \texttt{wsl\ -\/\allowbreak{}-set-version\ Ubuntu\ 2}. \texttt{*} өөр distribution дээр байвал: \texttt{wsl\ -\/\allowbreak{}-set-default\ Ubuntu}.

\hypertarget{ux430.1.4-git-for-windows}{%
\subsubsection{А.1.4 Git for Windows}\label{ux430.1.4-git-for-windows}}

\textbf{{[}ЗК · PowerShell{]}}

\begin{Shaded}
\begin{Highlighting}[]
\NormalTok{winget install {-}{-}id Git.}\FunctionTok{Git}\NormalTok{ {-}e}
\end{Highlighting}
\end{Shaded}

Суусны дараа \textbf{PowerShell-ээ хаагаад шинээр нээ} (шинэ програм PATH-д нэмэгдэнэ).

\textbf{{[}ЗК · PowerShell{]}}

\begin{Shaded}
\begin{Highlighting}[]
\NormalTok{git {-}{-}version}
\NormalTok{git credential{-}manager {-}{-}version}
\end{Highlighting}
\end{Shaded}

✓ \textbf{Шалгах:} хоёулаа хувилбарын дугаар хэвлэнэ (\texttt{git\ version\ 2.\allowbreak{}5x\ldots{}}, \texttt{2.x.x\ldots{}}).

\begin{quote}
☞ Git-ийг Ubuntu дотор ч ашиглана (А.1.9). Windows-ийн Git нь Ubuntu-д GitHub-ын нэвтрэлтийг хадгалах \textbf{Git Credential Manager}-ийг өгдөг (А.3).
\end{quote}

\hypertarget{ux430.1.5-vs-code-ux431ux430-ux4e9ux440ux433ux4e9ux442ux433ux4e9ux43bux4afux4afux434}{%
\subsubsection{А.1.5 VS Code ба өргөтгөлүүд}\label{ux430.1.5-vs-code-ux431ux430-ux4e9ux440ux433ux4e9ux442ux433ux4e9ux43bux4afux4afux434}}

\textbf{{[}ЗК · PowerShell{]}}

\begin{Shaded}
\begin{Highlighting}[]
\NormalTok{winget install {-}{-}id Microsoft.}\FunctionTok{VisualStudioCode}\NormalTok{ {-}e}
\end{Highlighting}
\end{Shaded}

PowerShell-ээ \textbf{хаагаад шинээр нээ}, дараа нь өргөтгөлүүдийг суулга (мөр бүрийг тус тусад нь, эсвэл бүгдийг нэг дор буулгаж болно):

\textbf{{[}ЗК · PowerShell{]}}

\begin{Shaded}
\begin{Highlighting}[]
\NormalTok{code {-}{-}install{-}extension ms{-}vscode{-}remote.}\FunctionTok{remote}\NormalTok{{-}wsl}
\NormalTok{code {-}{-}install{-}extension ms{-}vscode{-}remote.}\FunctionTok{remote}\NormalTok{{-}ssh}
\NormalTok{code {-}{-}install{-}extension ms{-}python.}\FunctionTok{python}
\NormalTok{code {-}{-}install{-}extension ms{-}azuretools.}\FunctionTok{vscode}\NormalTok{{-}containers}
\NormalTok{code {-}{-}install{-}extension redhat.}\FunctionTok{vscode}\NormalTok{{-}yaml}
\end{Highlighting}
\end{Shaded}

\begingroup\scriptsize\begin{longtable}[]{@{}ll@{}}
\toprule
Өргөтгөл & Юунд\tabularnewline
\midrule
\endhead
\textbf{WSL} & Ubuntu доторх файлыг VS Code-оор нээх (\texttt{code\ .})\tabularnewline
\textbf{Remote - SSH} & Pi дээрх файлыг шууд засах (В.3)\tabularnewline
\textbf{Python} & Python код\tabularnewline
\textbf{Container Tools} & Контейнер, Compose файл\tabularnewline
\textbf{YAML} & \texttt{docker-compose.\allowbreak{}yml}-ийн алдааг илрүүлэх\tabularnewline
\bottomrule
\end{longtable}\endgroup

✓ \textbf{Шалгах:} \texttt{code\ -\/\allowbreak{}-list-extensions} → дээрх таван ID жагсаалтад байна.

\hypertarget{ux430.1.6-docker-desktop}{%
\subsubsection{А.1.6 Docker Desktop}\label{ux430.1.6-docker-desktop}}

\textbf{Алхам 1 --- суулгах.}

\textbf{{[}ЗК · PowerShell{]}}

\begin{Shaded}
\begin{Highlighting}[]
\NormalTok{winget install {-}{-}id Docker.}\FunctionTok{DockerDesktop}\NormalTok{ {-}e}
\end{Highlighting}
\end{Shaded}

Installer-ийн цонх нээгдвэл:

\begin{enumerate}
\def\labelenumi{\arabic{enumi}.}
\tightlist
\item
  \textbf{Configuration} хуудсанд \textbf{Use WSL 2 instead of Hyper-V} чагттай байгааг шалга (анхдагчаар чагттай).
\item
  \textbf{OK} дар → суулгаж дуусахыг хүлээ (3--5 мин).
\item
  \textbf{Close} (эсвэл \textbf{Close and log out}) гарвал дар. Log out хийвэл дахин нэвтэр.
\end{enumerate}

\begin{quote}
☞ \texttt{winget} ажиллахгүй бол: https://docs.docker.com/desktop/setup/install/windows-install/ → \textbf{Docker Desktop for Windows -- x86\_64} товчоор \texttt{Docker\ Desktop\ Installer.\allowbreak{}exe} татаж, давхар товшоод дээрх 1--3-ыг хий.
\end{quote}

\textbf{Алхам 2 --- анх асаах.}

\begin{enumerate}
\def\labelenumi{\arabic{enumi}.}
\tightlist
\item
  Start → \texttt{Docker\ Desktop} → нээ.
\item
  \textbf{Docker Subscription Service Agreement} гарна → \textbf{Accept}. (Зөвшөөрөхгүй бол Docker Desktop ажиллахгүй. Боловсролын зорилгоор үнэгүй.)
\item
  Нэвтрэх (Sign in) эсвэл судалгааны цонх гарвал \textbf{Skip} / \textbf{Continue without signing in} дар.
\item
  Цонхны \textbf{зүүн доод буланд} \texttt{Engine\ running} (ногоон) гэж гартал хүлээ --- анх удаа 1--3 мин.
\end{enumerate}

\textbf{Алхам 3 --- WSL-тэй холбох.} Docker Desktop-ийн баруун дээд буланд \textbf{⚙ (Settings)}:

\begin{enumerate}
\def\labelenumi{\arabic{enumi}.}
\tightlist
\item
  \textbf{General} → \textbf{Use the WSL 2 based engine} чагттай байгааг шалга.
\item
  \textbf{Resources} → \textbf{WSL integration} →

  \begin{itemize}
  \tightlist
  \item
    \textbf{Enable integration with my default WSL distro} --- чагттай,
  \item
    доорх жагсаалтад \textbf{Ubuntu}-гийн шилжүүлэгч --- \textbf{асаалттай} (цэнхэр).
  \end{itemize}
\item
  \textbf{Apply \& restart} дар.
\end{enumerate}

\textbf{Алхам 4 --- шалгах (хоёр терминалаас).}

\textbf{{[}ЗК · PowerShell{]}}

\begin{Shaded}
\begin{Highlighting}[]
\NormalTok{docker version}
\NormalTok{docker compose version}
\NormalTok{docker run {-}{-}rm hello{-}world}
\end{Highlighting}
\end{Shaded}

\textbf{{[}ЗК · Ubuntu{]}}

\begin{Shaded}
\begin{Highlighting}[]
\ExtensionTok{docker}\NormalTok{ version}
\ExtensionTok{docker}\NormalTok{ compose version}
\ExtensionTok{docker}\NormalTok{ run {-}{-}rm hello{-}world}
\end{Highlighting}
\end{Shaded}

✓ \textbf{Шалгах (хоёр терминал дээр хоёуланд нь):}

\begin{itemize}
\tightlist
\item
  \texttt{docker\ version} нь \textbf{Client:} ба \textbf{Server:} гэсэн \textbf{хоёр} хэсэг хэвлэнэ.
\item
  \texttt{docker\ compose\ version} → \texttt{Docker\ Compose\ version\ v2.\allowbreak{}x.\allowbreak{}x}.
\item
  \texttt{hello-world} → \texttt{Hello\ from\ Docker!} гэсэн мөр гарна.
\end{itemize}

✗ \texttt{error\ during\ connect} / \texttt{Cannot\ connect\ to\ the\ Docker\ daemon} → Docker Desktop ажиллаагүй байна: Алхам 2-ын 4-р зүйлийг шалга.
✗ Ubuntu-д \texttt{docker:\allowbreak{}\ command\ not\ found} эсвэл \texttt{The\ command\ \textquotesingle{}docker\textquotesingle{}\ could\ not\ be\ found\ in\ this\ WSL\ 2\ distro} → Алхам 3-ын \textbf{WSL integration} асаагүй байна.

\begin{quote}
☞ \textbf{Docker Desktop-ийг лаб бүрийн өмнө асаа.} Компьютер асах бүрт автоматаар асаах бол: \textbf{Settings → General → Start Docker Desktop when you sign in to your computer}.
\end{quote}

\hypertarget{ux430.1.7-raspberry-pi-imager}{%
\subsubsection{А.1.7 Raspberry Pi Imager}\label{ux430.1.7-raspberry-pi-imager}}

\textbf{{[}ЗК · PowerShell{]}}

\begin{Shaded}
\begin{Highlighting}[]
\NormalTok{winget install {-}{-}id RaspberryPiFoundation.}\FunctionTok{RaspberryPiImager}\NormalTok{ {-}e}
\end{Highlighting}
\end{Shaded}

✓ \textbf{Шалгах:} Start цэсэнд \textbf{Raspberry Pi Imager} гарч ирнэ. Хувилбар нь \textbf{2.0 ба түүнээс шинэ} байх ёстой (Б.1-ийн дэлгэцүүд энэ хувилбарынх).

\begin{quote}
☞ Гараар: https://www.raspberrypi.com/software/ → \textbf{Download for Windows}.
\end{quote}

\hypertarget{ux430.1.8-mqtt-explorer-ux441ux43eux43dux433ux43eux43bux442ux43eux442}{%
\subsubsection{А.1.8 MQTT Explorer (сонголтот)}\label{ux430.1.8-mqtt-explorer-ux441ux43eux43dux433ux43eux43bux442ux43eux442}}

Брокер дахь сэдвийн модыг график хэлбэрээр харуулна. Заавал биш --- \texttt{mosquitto\_sub} хангалттай.

\textbf{{[}ЗК · PowerShell{]}}

\begin{Shaded}
\begin{Highlighting}[]
\NormalTok{winget install {-}{-}id thomasnordquist.}\FunctionTok{MQTT}\NormalTok{{-}Explorer {-}e}
\end{Highlighting}
\end{Shaded}

\hypertarget{ux430.1.9-ubuntu-ux434ux43eux442ux43eux440ux445-ux445ux44dux440ux44dux433ux441ux43bux4afux4afux434}{%
\subsubsection{А.1.9 Ubuntu доторх хэрэгслүүд}\label{ux430.1.9-ubuntu-ux434ux43eux442ux43eux440ux445-ux445ux44dux440ux44dux433ux441ux43bux4afux4afux434}}

\textbf{{[}ЗК · Ubuntu{]}}

\begin{Shaded}
\begin{Highlighting}[]
\FunctionTok{sudo}\NormalTok{ apt update}
\FunctionTok{sudo}\NormalTok{ apt upgrade {-}y}
\FunctionTok{sudo}\NormalTok{ apt install {-}y git make curl bc jq python3{-}venv python3{-}pip mosquitto{-}clients}
\end{Highlighting}
\end{Shaded}

Эхний команд Ubuntu-гийн нууц үгийг (А.1.3, Алхам 2) асууна.

✓ \textbf{Шалгах:}

\textbf{{[}ЗК · Ubuntu{]}}

\begin{Shaded}
\begin{Highlighting}[]
\FunctionTok{git}\NormalTok{ {-}{-}version}\KeywordTok{;} \FunctionTok{make}\NormalTok{ {-}{-}version }\KeywordTok{|} \FunctionTok{head}\NormalTok{ {-}1}\KeywordTok{;} \ExtensionTok{python3}\NormalTok{ {-}{-}version}\KeywordTok{;} \ExtensionTok{mosquitto\_sub}\NormalTok{ {-}{-}help }\KeywordTok{|} \FunctionTok{head}\NormalTok{ {-}1}
\end{Highlighting}
\end{Shaded}

Дөрвөн мөр хувилбарын мэдээлэл хэвлэнэ (\texttt{python3} нь \textbf{3.10+}).

\hypertarget{ux430.2-docker-desktop-ux438ux439ux43d-ux441ux430ux43dux430ux445-ux43eux439-ux437ux430ux430ux432ux430ux43b}{%
\subsection{А.2 Docker Desktop-ийн санах ой (ЗААВАЛ)}\label{ux430.2-docker-desktop-ux438ux439ux43d-ux441ux430ux43dux430ux445-ux43eux439-ux437ux430ux430ux432ux430ux43b}}

Docker Desktop-ийн баримтаар Linux VM-д анхдагчаар \textbf{хост компьютерийн санах ойн 50\%}-ийг өгдөг. Бидэнд \textbf{6 GB} хэрэгтэй (Лаб 8-д Ollama). Тохируулах газар нь backend-ээс хамаарна:

\begin{itemize}
\tightlist
\item
  \textbf{Windows + WSL 2 backend (бидний сонголт):} Docker Desktop-ийн \textbf{Resources} хэсэгт Memory гулсуур \textbf{байхгүй} --- санах ойг WSL 2-ын VM-д \texttt{\%UserProfile\%\textbackslash{}.wslconfig} файлаар өгнө (доорх алхмууд).
\item
  \textbf{macOS, Linux, Windows Hyper-V backend:} \textbf{Settings → Resources → Advanced → Memory limit → 6 GB} (16 GB-тай бол 8 GB) → \emph{Apply \& restart}.
\end{itemize}

\textbf{Алхам 1 --- файл нээх.}

\textbf{{[}ЗК · PowerShell{]}}

\begin{Shaded}
\begin{Highlighting}[]
\NormalTok{notepad }\StringTok{"$env:USERPROFILE\textbackslash{}.wslconfig"}
\end{Highlighting}
\end{Shaded}

\texttt{Cannot\ find\ the\ .\allowbreak{}wslconfig\ file.\allowbreak{}\ Do\ you\ want\ to\ create\ a\ new\ file?\allowbreak{}} гэж асуувал \textbf{Yes}.

\textbf{Алхам 2 --- дараах хоёр мөрийг бичээд хадгал} (\texttt{Ctrl\ +\ S}), Notepad-ийг хаа:

\begin{Shaded}
\begin{Highlighting}[]
\KeywordTok{[wsl2]}
\DataTypeTok{memory}\OtherTok{=}\StringTok{6GB}
\end{Highlighting}
\end{Shaded}

\begingroup\scriptsize\begin{longtable}[]{@{}ll@{}}
\toprule
Компьютерийн RAM & \texttt{memory=}\tabularnewline
\midrule
\endhead
8 GB & \texttt{6GB} (Лаб 8-д Chrome ба бусад програмаа хаана)\tabularnewline
16 GB ба түүнээс их & \texttt{8GB}\tabularnewline
\bottomrule
\end{longtable}\endgroup

\begin{quote}
⚠ Файлын нэр яг \texttt{.wslconfig} байх ёстой --- \textbf{\texttt{.wslconfig.txt} БИШ}. Шалгах: \texttt{Get-ChildItem\ \$env:USERPROFILE\textbackslash{}.wslconfig} → алдаагүй гарвал зөв.
\end{quote}

\textbf{Алхам 3 --- WSL-ийг дахин эхлүүлэх.}

\begin{enumerate}
\def\labelenumi{\arabic{enumi}.}
\tightlist
\item
  Docker Desktop-ийг хаа: taskbar-ын баруун доод талын \textbf{халимны дүрс} дээр баруун товч → \textbf{Quit Docker Desktop}.
\item
  Бүх Ubuntu цонхоо хаа.
\item
  \textbf{{[}ЗК · PowerShell{]}} \texttt{wsl\ -\/\allowbreak{}-shutdown}
\item
  10 секунд хүлээгээд Docker Desktop-ийг дахин нээ → \texttt{Engine\ running} болтол хүлээ.
\end{enumerate}

✓ \textbf{Шалгах:}

\textbf{{[}ЗК · Ubuntu{]}}

\begin{Shaded}
\begin{Highlighting}[]
\ExtensionTok{docker}\NormalTok{ info {-}{-}format }\StringTok{\textquotesingle{}\{\{.MemTotal\}\}\textquotesingle{}} \KeywordTok{|} \FunctionTok{awk} \StringTok{\textquotesingle{}\{printf "Docker{-}т \%.1f GB\textbackslash{}n", $1/1024/1024/1024\}\textquotesingle{}}
\end{Highlighting}
\end{Shaded}

\texttt{Docker-т\ 5.\allowbreak{}8\ GB} орчим (6GB тавьсан бол 5.7--6.0) гарна. Хэвээр 50\% бол файлын нэр эсвэл байршил буруу (Алхам 2-ын ⚠). \textbf{4 GB-аас бага бол Лаб 8-д \texttt{ai} профайл ажиллахгүй.}

\begin{quote}
Албан ёсны баримт: \href{https://docs.docker.com/desktop/settings-and-maintenance/settings/}{Docker Desktop --- Settings (Resources)} · \href{https://learn.microsoft.com/en-us/windows/wsl/wsl-config}{WSL --- .wslconfig}
\end{quote}

\hypertarget{ux430.3-ux441ux430ux43dux433ux430ux430-ux442ux430ux442ux430ux445-ux431ux430-python-ux43eux440ux447ux438ux43d}{%
\subsection{А.3 Сангаа татах ба Python орчин}\label{ux430.3-ux441ux430ux43dux433ux430ux430-ux442ux430ux442ux430ux445-ux431ux430-python-ux43eux440ux447ux438ux43d}}

\textbf{Алхам 1 --- Git тохируулах} (нэг удаа):

\textbf{{[}ЗК · Ubuntu{]}}

\begin{Shaded}
\begin{Highlighting}[]
\FunctionTok{git}\NormalTok{ config {-}{-}global user.name }\StringTok{"Овог Нэр"}
\FunctionTok{git}\NormalTok{ config {-}{-}global user.email }\StringTok{"таны@имэйл.mn"}
\FunctionTok{git}\NormalTok{ config {-}{-}global core.autocrlf input}
\FunctionTok{git}\NormalTok{ config {-}{-}global credential.helper }\StringTok{"/mnt/c/Program\textbackslash{} Files/Git/mingw64/bin/git{-}credential{-}manager.exe"}
\end{Highlighting}
\end{Shaded}

\begin{itemize}
\tightlist
\item
  \texttt{core.\allowbreak{}autocrlf\ input} --- мөрийн төгсгөлийг Linux-ийн (LF) хэвээр үлдээнэ. Үгүй бол \texttt{.sh} файлд CRLF орж контейнерт \texttt{/bin/sh\^{}M:\ bad\ interpreter} алдаа гарна.
\item
  \texttt{credential.\allowbreak{}helper} --- GitHub-ын нэвтрэлтийг Windows-ийн Git Credential Manager-т хадгална: \texttt{git\ clone}/\allowbreak\texttt{git\ push} анх удаа \textbf{хөтөч нээж} нэвтрүүлнэ, дараа нь дахин асуухгүй.
\end{itemize}

\textbf{Алхам 2 --- сангаа Ubuntu-гийн гэрийн хавтаст clone хийх.}

\begin{quote}
⚠ Сангаа \textbf{Ubuntu дотор} (\texttt{\textasciitilde{}/cnc302}) хадгална --- \texttt{/mnt/c/\ldots{}}, \textbf{OneDrive}, \textbf{Desktop}, \textbf{Documents}-д БИШ. Docker-ийн WSL-ийн зөвлөмжөөр Linux файлын системд байгаа файл контейнерт хамаагүй хурдан харагдана; OneDrive нь \texttt{.git}-ийг эвддэг.
\end{quote}

\textbf{{[}ЗК · Ubuntu{]}}

\begin{Shaded}
\begin{Highlighting}[]
\BuiltInTok{cd}\NormalTok{ \textasciitilde{}}
\FunctionTok{git}\NormalTok{ clone }\OperatorTok{\textless{}}\NormalTok{REPO\_URL}\OperatorTok{\textgreater{}}\NormalTok{ cnc302}
\BuiltInTok{cd}\NormalTok{ \textasciitilde{}/cnc302}
\FunctionTok{ls}
\end{Highlighting}
\end{Shaded}

✓ \textbf{Шалгах:} \texttt{lab01\ \ldots{}\ lab08\ \ stack\ \ edge\ \ tools\ \ docs\ \ README.\allowbreak{}md\ \ SETUP.\allowbreak{}md} харагдана.

\begin{quote}
☞ Сангаа VS Code-оор нээх: \texttt{cd\ \textasciitilde{}/cnc302\ \&\&\ code\ .} --- VS Code нээгдэж, зүүн доод буланд \textbf{WSL: Ubuntu} гэж харагдана. Windows-ийн File Explorer-оос: хаягийн мөрөнд \texttt{\textbackslash{}\textbackslash{}wsl\$\textbackslash{}Ubuntu\textbackslash{}home\textbackslash{}\textless{}ubuntu-нэр\textgreater{}\textbackslash{}cnc302}.
\end{quote}

\textbf{Алхам 3 --- Python-ийн virtual environment.}

\textbf{{[}ЗК · Ubuntu{]}}

\begin{Shaded}
\begin{Highlighting}[]
\BuiltInTok{cd}\NormalTok{ \textasciitilde{}/cnc302}
\ExtensionTok{python3}\NormalTok{ {-}m venv .venv}
\BuiltInTok{source}\NormalTok{ .venv/bin/activate}
\ExtensionTok{pip}\NormalTok{ install {-}{-}upgrade pip}
\ExtensionTok{pip}\NormalTok{ install {-}r tools/requirements.txt}
\end{Highlighting}
\end{Shaded}

✓ Prompt-ын эхэнд \texttt{(.venv)} гарна. Шинэ Ubuntu цонх нээх бүрт \textbf{\texttt{cd\ \textasciitilde{}/cnc302\ \&\&\ source\ .venv/bin/activate}}-ийг дахин бичнэ.

\textbf{Алхам 4 --- шалгах.}

\textbf{{[}ЗК · Ubuntu{]}}

\begin{Shaded}
\begin{Highlighting}[]
\ExtensionTok{python}\NormalTok{ tools/sim\_device.py {-}{-}dry{-}run {-}{-}devices 2 {-}{-}count 3}
\end{Highlighting}
\end{Shaded}

✓ Брокергүйгээр \textbf{6 мессежийн JSON} хэвлэгдэнэ. Хэвлэгдэхгүй бол цааш явахгүй.

\hypertarget{ux430.4-ux4afux4afux43bux43dux438ux439-ux441ux442ux435ux43aux438ux439ux433-ux442ux430ux442ux430ux436-ux430ux441ux430ux430ux445}{%
\subsection{А.4 Үүлний стекийг татаж, асаах}\label{ux430.4-ux4afux4afux43bux43dux438ux439-ux441ux442ux435ux43aux438ux439ux433-ux442ux430ux442ux430ux436-ux430ux441ux430ux430ux445}}

\textbf{Алхам 1 --- нууц утгын файл үүсгэх.}

\textbf{{[}ЗК · Ubuntu{]}}

\begin{Shaded}
\begin{Highlighting}[]
\BuiltInTok{cd}\NormalTok{ \textasciitilde{}/cnc302/stack}
\FunctionTok{cp}\NormalTok{ .env.example .env}
\FunctionTok{nano}\NormalTok{ .env}
\end{Highlighting}
\end{Shaded}

\texttt{nano} засварлагч нээгдэнэ. Сумаар зөөж дараах \textbf{гурван утгыг} өөрийнхөөрөө соль:

\begingroup\scriptsize\begin{longtable}[]{@{}ll@{}}
\toprule
Мөр & Юу бичих\tabularnewline
\midrule
\endhead
\texttt{EMQX\_COOKIE=} & дурын урт үг (жишээ: \texttt{team07-cookie-x8k2})\tabularnewline
\texttt{EMQX\_\allowbreak{}DASHBOARD\_\allowbreak{}PASSWORD=\allowbreak{}} & EMQX самбарын нууц үг --- \textbf{тэмдэглэж ав}\tabularnewline
\texttt{GRAFANA\_\allowbreak{}PASSWORD=\allowbreak{}} & Grafana-гийн нууц үг --- \textbf{тэмдэглэж ав}\tabularnewline
\bottomrule
\end{longtable}\endgroup

Бусад мөрийг (жишээ нь хоосон \texttt{EMQX\_\allowbreak{}API\_\allowbreak{}KEY=\allowbreak{}}) \textbf{хэвээр нь} үлдээ --- тэдгээрийг Лаб 2-т бөглөнө.

Хадгалах: \texttt{Ctrl\ +\ O} → \texttt{Enter}. Гарах: \texttt{Ctrl\ +\ X}.

\begin{quote}
⚠ Нууц үгийг \textbf{стекийг анх асаахаас өмнө} соль. EMQX ба Grafana нууц үгээ зөвхөн \textbf{анх асахдаа} \texttt{.env}-ээс уншина; дараа нь \texttt{.env}-ийг сольсон ч өөрчлөгдөхгүй.
\end{quote}

\textbf{Алхам 2 --- image-уудыг татах} (анх удаа хэдэн GB, \textbf{10--30 мин}):

\textbf{{[}ЗК · Ubuntu{]}}

\begin{Shaded}
\begin{Highlighting}[]
\ExtensionTok{docker}\NormalTok{ compose {-}{-}profile core pull {-}{-}ignore{-}buildable}
\end{Highlighting}
\end{Shaded}

\texttt{-\/\allowbreak{}-ignore-buildable} --- манай \texttt{registry} image Docker Hub-д байхгүй, дараагийн алхамд \textbf{локалд барина}; энэ туг түүнийг татах гэж оролдохгүй болгоно.

\textbf{Алхам 3 --- асаах.}

\textbf{{[}ЗК · Ubuntu{]}}

\begin{Shaded}
\begin{Highlighting}[]
\FunctionTok{make}\NormalTok{ up}
\end{Highlighting}
\end{Shaded}

Анх удаа \texttt{registry}-г барина (2--5 мин). Төгсгөлд \texttt{Pi-гийн\ edge/\allowbreak{}.\allowbreak{}env\ файлд\ бичих\ утга:\allowbreak{}} гэсэн жагсаалт хэвлэгдэнэ --- А.5-д хэрэглэнэ.

\textbf{Алхам 4 --- шалгах.} 1 минут хүлээгээд:

\textbf{{[}ЗК · Ubuntu{]}}

\begin{Shaded}
\begin{Highlighting}[]
\FunctionTok{make}\NormalTok{ status}
\FunctionTok{make}\NormalTok{ health}
\end{Highlighting}
\end{Shaded}

✓ \textbf{Шалгах:}

\begin{itemize}
\item
  \texttt{make\ status}-ын \textbf{STATUS} баганад бүх контейнер \texttt{Up} эсвэл \texttt{Up\ (healthy)}. \texttt{starting} / \texttt{health:\allowbreak{}\ starting} бол 1 мин хүлээгээд дахин ажиллуул.
\item
  \texttt{make\ health} дөрвөн мөр \textbf{\texttt{200}} хэвлэнэ:

\begin{verbatim}
EMQX      200
InfluxDB  200
Grafana   200
Registry  200
\end{verbatim}
\end{itemize}

Хөтчөөр шалгах:

\begingroup\scriptsize\begin{longtable}[]{@{}ll@{}}
\toprule
Хаяг & Нэвтрэх\tabularnewline
\midrule
\endhead
http://localhost:18083 & EMQX самбар: \texttt{admin} / \texttt{EMQX\_\allowbreak{}DASHBOARD\_\allowbreak{}PASSWORD}\tabularnewline
http://localhost:3000 & Grafana: \texttt{admin} / \texttt{GRAFANA\_\allowbreak{}PASSWORD}\tabularnewline
\bottomrule
\end{longtable}\endgroup

\begin{quote}
☞ Стекийг зогсоох (өгөгдөл устахгүй): \texttt{make\ down}. Дахин асаах: \texttt{make\ up}. \textbf{\texttt{make\ clean} бүх өгөгдлийг устгана} --- хэрэглэхээсээ өмнө бод.
\end{quote}

\hypertarget{ux430.5-lan-ip-ux445ux430ux44fux433ux430ux430-ux43eux43bux43eux445-ux447ux443ux445ux430ux43b}{%
\subsection{А.5 LAN IP хаягаа олох (ЧУХАЛ)}\label{ux430.5-lan-ip-ux445ux430ux44fux433ux430ux430-ux43eux43bux43eux445-ux447ux443ux445ux430ux43b}}

Pi нь зөөврийн компьютерт \textbf{LAN IP}-ээр холбогдоно --- \texttt{localhost} БИШ (Pi дээр \texttt{localhost} гэдэг нь Pi өөрөө).

\textbf{{[}ЗК · PowerShell{]}}

\begin{Shaded}
\begin{Highlighting}[]
\NormalTok{ipconfig}
\end{Highlighting}
\end{Shaded}

Гаралтаас \textbf{Wireless LAN adapter Wi-Fi} (кабелиар бол \textbf{Ethernet adapter Ethernet}) хэсгийг олж, \textbf{IPv4 Address} мөрийн утгыг ав. Энэ бол \texttt{\textless{}LAPTOP\_\allowbreak{}IP\textgreater{}} --- \textbf{тэмдэглэ}.

\begin{quote}
⚠ \texttt{vEthernet\ (WSL\ldots{})} эсвэл \texttt{vEthernet\ (Default\ Switch)} хэсгийн хаягийг БҮҮ ав --- тэр нь Windows доторх \textbf{дотоод} сүлжээ, Pi түүнийг харахгүй. Ubuntu дотор \texttt{ip\ addr}-оор гарах \texttt{172.x.x.x} хаяг ч мөн адил \textbf{буруу}.

☞ Ubuntu-гийн \texttt{make\ ip} (\texttt{make\ up}-ын төгсгөлд хэвлэгддэг) WSL дотор Windows-ийн \texttt{ipconfig}-ийн IPv4 мөрүүдийг харуулна --- тэндээс мөн Wi-Fi/Ethernet адаптерийнхийг сонго.
\end{quote}

\textbf{Сүлжээний профайлыг Private болгох.} \textbf{Settings → Network \& internet → Wi-Fi} (эсвэл \textbf{Ethernet}) → холбогдсон сүлжээний нэр → \textbf{Network profile type} → \textbf{Private network}. Public байвал Windows Firewall Pi-гаас ирэх холболтыг хаана (А.6).

\begin{quote}
\textbf{Сүлжээний урхи:}
- Хамгийн найдвартай нь \textbf{хоёуланг нь нэг router/switch-д кабелиар} холбох.
- Pi 3B-ийн Wi-Fi нь \textbf{зөвхөн 2.4 GHz}. 5 GHz сүлжээ Imager-т харагдах боловч Pi холбогдохгүй.
- Их сургуулийн нийтийн Wi-Fi ихэвчлэн төхөөрөмжүүдийг хооронд нь \textbf{тусгаарладаг} (client isolation) тул Pi зөөврийн компьютерээ "харахгүй". Лабораторийн router эсвэл утасныхаа \textbf{2.4 GHz hotspot}-ыг ашигла.
- Зөөврийн компьютер Wi-Fi-аар, Pi кабелиар холбогдсон бол өөр дэд сүлжээнд (subnet) байж болно --- \texttt{\textless{}LAPTOP\_\allowbreak{}IP\textgreater{}} ба \texttt{\textless{}PI\_\allowbreak{}IP\textgreater{}}-ийн эхний гурван тоо ижил эсэхийг шалга.
- IP \textbf{өөрчлөгдөж болно} (өөр сүлжээ, router restart). Лаб бүрийн эхэнд \texttt{ipconfig}-оор шалга.
\end{quote}

\hypertarget{ux430.6-ux433ux430ux43bux442-ux445ux430ux43dux430-windows}{%
\subsection{А.6 Галт хана (Windows)}\label{ux430.6-ux433ux430ux43bux442-ux445ux430ux43dux430-windows}}

Windows Defender Firewall гаднаас ирэх холболтыг анхдагчаар хаадаг. Pi-гаас ирэх портуудыг нээнэ: \textbf{1883} (MQTT, гүүр), \textbf{8883} (MQTT/TLS, Лаб 2), \textbf{8090} (registry/OTA, Лаб 2), \textbf{8181} (InfluxDB), \textbf{3000} (Grafana).

\textbf{{[}ЗК · PowerShell (Admin){]}} --- доорх \textbf{хоёр мөрийг бүтнээр нь нэг дор} хуулж буулга (эхний мөрийн төгсгөлийн \texttt{\textasciigrave{}} тэмдэг нь "команд дараагийн мөрөнд үргэлжилнэ" гэсэн утгатай):

\begin{Shaded}
\begin{Highlighting}[]
\NormalTok{New{-}NetFirewallRule {-}DisplayName }\StringTok{"CNC302 cloud"}\NormalTok{ {-}Direction Inbound \textasciigrave{}}
\NormalTok{  {-}Protocol TCP {-}LocalPort 1883,8883,8090,8181,3000 {-}Action Allow {-}Profile Private}
\end{Highlighting}
\end{Shaded}

✓ \textbf{Шалгах:}

\textbf{{[}ЗК · PowerShell{]}}

\begin{Shaded}
\begin{Highlighting}[]
\NormalTok{Get{-}NetFirewallRule {-}DisplayName }\StringTok{"CNC302 cloud"}\NormalTok{ | }\FunctionTok{Format{-}Table}\NormalTok{ DisplayName, Enabled, Profile, Action}
\end{Highlighting}
\end{Shaded}

\texttt{CNC302\ cloud\ \ True\ \ Private\ \ Allow} гарна.

\begin{quote}
Зөвхөн \textbf{Private} профайлд нээ (А.5). Нийтийн сүлжээнд задгай MQTT брокер орхиж болохгүй.

☞ Docker Desktop эсвэл өөр програм \textbf{"Windows Defender Firewall has blocked some features"} цонх гаргавал \textbf{Private networks}-ийг чагтлаад \textbf{Allow access} дар.

Албан ёсны баримт: \href{https://learn.microsoft.com/en-us/powershell/module/netsecurity/new-netfirewallrule}{New-NetFirewallRule}
\end{quote}

\hypertarget{ux430.7-ssh-ux442ux4afux43bux445ux4afux4afux440-ux4afux4afux441ux433ux44dux445}{%
\subsection{А.7 SSH түлхүүр үүсгэх}\label{ux430.7-ssh-ux442ux4afux43bux445ux4afux4afux440-ux4afux4afux441ux433ux44dux445}}

Pi руу \textbf{нууц үггүй, түлхүүрээр} нэвтэрнэ. Түлхүүрийг \textbf{Windows} талд үүсгэнэ --- PowerShell-ийн \texttt{ssh} болон VS Code хоёулаа үүнийг ашиглана.

\textbf{Алхам 1 --- түлхүүр үүсгэх.}

\textbf{{[}ЗК · PowerShell{]}}

\begin{Shaded}
\begin{Highlighting}[]
\FunctionTok{New{-}Item}\NormalTok{ {-}ItemType Directory {-}Force }\StringTok{"$env:USERPROFILE\textbackslash{}.ssh"}\NormalTok{ | }\FunctionTok{Out{-}Null}
\NormalTok{ssh{-}keygen {-}t ed25519 {-}C }\StringTok{"таны@имэйл.mn"}\NormalTok{ {-}f }\StringTok{"$env:USERPROFILE\textbackslash{}.ssh\textbackslash{}cnc302"}
\end{Highlighting}
\end{Shaded}

\texttt{Enter\ passphrase} гэж асуувал хоосон орхиж \textbf{Enter}, дахин \textbf{Enter}.

✓ \texttt{Your\ identification\ has\ been\ saved\ in\ C:\textbackslash{}Users\textbackslash{}\ldots{}\textbackslash{}.ssh\textbackslash{}cnc302} гэж гарна. Хоёр файл үүснэ:

\begingroup\scriptsize\begin{longtable}[]{@{}lll@{}}
\toprule
Файл & Юу вэ &\tabularnewline
\midrule
\endhead
\texttt{cnc302} & \textbf{нууц} түлхүүр & СТОП: \textbf{Хэнд ч бүү өг, Git-д хэзээ ч бүү оруул}\tabularnewline
\texttt{cnc302.pub} & \textbf{нийтийн} түлхүүр & Pi-д өгнө (Б.1)\tabularnewline
\bottomrule
\end{longtable}\endgroup

\textbf{Алхам 2 --- нийтийн түлхүүрийг харах.} Б.1-д хуулж буулгана:

\textbf{{[}ЗК · PowerShell{]}}

\begin{Shaded}
\begin{Highlighting}[]
\FunctionTok{Get{-}Content} \StringTok{"$env:USERPROFILE\textbackslash{}.ssh\textbackslash{}cnc302.pub"}
\end{Highlighting}
\end{Shaded}

\texttt{ssh-ed25519\ AAAAC3Nza\ldots{}\ таны@имэйл.\allowbreak{}mn} гэсэн \textbf{нэг урт мөр} гарна.

\textbf{Алхам 3 --- \texttt{ssh\ pi} товчлол үүсгэх.} Доорх блокийн \texttt{\textless{}NN\textgreater{}}-ийг багийн дугаараар \textbf{сольсны дараа} бүтнээр нь буулга:

\textbf{{[}ЗК · PowerShell{]}}

\begin{Shaded}
\begin{Highlighting}[]
\FunctionTok{Add{-}Content}\NormalTok{ {-}Encoding ascii }\StringTok{"$env:USERPROFILE\textbackslash{}.ssh\textbackslash{}config"} \VerbatimStringTok{@"}

\VerbatimStringTok{Host pi}
\VerbatimStringTok{    HostName pi{-}team\textless{}NN\textgreater{}.local}
\VerbatimStringTok{    User cnc302}
\VerbatimStringTok{    IdentityFile \textasciitilde{}/.ssh/cnc302}
\VerbatimStringTok{"@}
\end{Highlighting}
\end{Shaded}

Одооноос \texttt{ssh\ pi} гэж бичихэд \texttt{ssh\ -i\ \textasciitilde{}/.ssh/cnc302\ cnc302@pi-team\textless{}NN\textgreater{}.local}-тай ижил ажиллана. VS Code ч энэ файлыг уншина (В.3).

✓ \textbf{Шалгах:} \texttt{Get-Content\ "\$env:USERPROFILE\textbackslash{}.ssh\textbackslash{}config"} → \texttt{HostName\ pi-team07.\allowbreak{}local} гэх мэт \textbf{өөрийн дугаартай} мөр харагдана (\texttt{\textless{}NN\textgreater{}} үлдсэн бол файлыг \texttt{notepad\ "\$env:USERPROFILE\textbackslash{}.ssh\textbackslash{}config"}-оор засна).

\hypertarget{ux430.8-vm-ux43bux430ux431-8-ux44bux43d-k3s-server-vm-ux43bux430ux431-8-ux430ux430ux441-ux4e9ux43cux43dux4e9-ux445ux438ux439ux43dux44d}{%
\subsection{А.8 {[}VM{]} Лаб 8-ын K3s server VM (Лаб 8-аас өмнө хийнэ)}\label{ux430.8-vm-ux43bux430ux431-8-ux44bux43d-k3s-server-vm-ux43bux430ux431-8-ux430ux430ux441-ux4e9ux43cux43dux4e9-ux445ux438ux439ux43dux44d}}

Лаб 8-д K3s-ийн \textbf{server} нь зөөврийн компьютер дээрх VM-д, Pi 3B нь \textbf{agent} болж нэгдэнэ. K3s-ийн албан ёсны доод шаардлага: server \textbf{2 цөм / 2 GB}, agent \textbf{1 цөм / 512 MB} --- 1 GB-тай Pi server-т хүрэлцэхгүй.

\begingroup\scriptsize\begin{longtable}[]{@{}ll@{}}
\toprule
\begin{minipage}[b]{0.47\columnwidth}\raggedright
Тохиргоо\strut
\end{minipage} & \begin{minipage}[b]{0.47\columnwidth}\raggedright
Утга\strut
\end{minipage}\tabularnewline
\midrule
\endhead
\begin{minipage}[t]{0.47\columnwidth}\raggedright
Гипервизор\strut
\end{minipage} & \begin{minipage}[t]{0.47\columnwidth}\raggedright
Oracle VirtualBox\strut
\end{minipage}\tabularnewline
\begin{minipage}[t]{0.47\columnwidth}\raggedright
Зочин OS\strut
\end{minipage} & \begin{minipage}[t]{0.47\columnwidth}\raggedright
\textbf{Ubuntu Server LTS} amd64 --- одоогийн LTS нь \textbf{26.04} (24.04 LTS ч дэмжигдсэн хэвээр)\strut
\end{minipage}\tabularnewline
\begin{minipage}[t]{0.47\columnwidth}\raggedright
vCPU / RAM / диск\strut
\end{minipage} & \begin{minipage}[t]{0.47\columnwidth}\raggedright
\textbf{≥ 2 vCPU}, \textbf{4 GB} (доод 2 GB), 20 GB\strut
\end{minipage}\tabularnewline
\begin{minipage}[t]{0.47\columnwidth}\raggedright
Сүлжээ\strut
\end{minipage} & \begin{minipage}[t]{0.47\columnwidth}\raggedright
\textbf{Bridged Adapter} (LAN-д холбогдсон физик адаптер), NAT адаптер нэмэхгүй\strut
\end{minipage}\tabularnewline
\begin{minipage}[t]{0.47\columnwidth}\raggedright
Суулгахдаа\strut
\end{minipage} & \begin{minipage}[t]{0.47\columnwidth}\raggedright
\textbf{OpenSSH server}-ийг сонго\strut
\end{minipage}\tabularnewline
\bottomrule
\end{longtable}\endgroup

Дэлгэрэнгүй алхам ба K3s-ийн суулгалт: \texttt{lab08/\allowbreak{}k3s/\allowbreak{}README.\allowbreak{}md} §3.1--3.2.

\begin{quote}
⚠ \textbf{Windows дээр Hyper-V ба VirtualBox:} Docker Desktop-ийн WSL 2 backend нь Windows-ийн гипервизорыг (Hyper-V платформ) асаадаг. VirtualBox-ийн гарын авлагаар Hyper-V ажиллаж байгаа хост дээр VirtualBox ажиллах боловч \emph{"significant Oracle VirtualBox performance degradation"} гарч болно; VirtualBox Hyper-V-г виртуалчлалын хөдөлгүүр болгон автоматаар ашиглах бөгөөд үүнд \textbf{Windows Hypervisor Platform}-ийг нэмж идэвхжүүлсэн байх шаардлагатай. VM-ийн цонхны CPU дүрс (яст мэлхий) энэ горимыг заана. VM-д 2 vCPU-гээс бага бүү өг.

Албан ёсны баримт: \href{https://www.virtualbox.org/manual/UserManual.html}{VirtualBox User Manual --- Using Hyper-V with Oracle VirtualBox} · \href{https://ubuntu.com/about/release-cycle}{Ubuntu release cycle} · \href{https://docs.k3s.io/installation/requirements}{K3s Requirements}
\end{quote}

\hypertarget{ux431.-pi-raspberry-pi-3b-ux434ux44dux44dux440}{%
\section{Б. {[}Pi{]} Raspberry Pi 3B дээр}\label{ux431.-pi-raspberry-pi-3b-ux434ux44dux44dux440}}

\hypertarget{ux431.0-ux442ux43eux43dux43eux433-ux442ux4e9ux445ux4e9ux4e9ux440ux4e9ux43cux436ux438ux439ux43d-ux448ux430ux43bux433ux443ux443ux440}{%
\subsection{Б.0 Тоног төхөөрөмжийн шалгуур}\label{ux431.0-ux442ux43eux43dux43eux433-ux442ux4e9ux445ux4e9ux4e9ux440ux4e9ux43cux436ux438ux439ux43d-ux448ux430ux43bux433ux443ux443ux440}}

\begingroup\scriptsize\begin{longtable}[]{@{}lll@{}}
\toprule
\begin{minipage}[b]{0.30\columnwidth}\raggedright
Зүйл\strut
\end{minipage} & \begin{minipage}[b]{0.30\columnwidth}\raggedright
Шаардлага\strut
\end{minipage} & \begin{minipage}[b]{0.30\columnwidth}\raggedright
Яагаад\strut
\end{minipage}\tabularnewline
\midrule
\endhead
\begin{minipage}[t]{0.30\columnwidth}\raggedright
Тэжээл\strut
\end{minipage} & \begin{minipage}[t]{0.30\columnwidth}\raggedright
\textbf{5 V / 2.5 A} (албан ёсны 12.5 W Micro USB адаптер)\strut
\end{minipage} & \begin{minipage}[t]{0.30\columnwidth}\raggedright
Сул тэжээл → хүчдэл унана (\texttt{get\_throttled}-ийн бит 0) → бүх хэмжилт гажина\strut
\end{minipage}\tabularnewline
\begin{minipage}[t]{0.30\columnwidth}\raggedright
Хөргөлт\strut
\end{minipage} & \begin{minipage}[t]{0.30\columnwidth}\raggedright
Наалдацтай хөргөгч, боломжтой бол сэнс\strut
\end{minipage} & \begin{minipage}[t]{0.30\columnwidth}\raggedright
Албан ёсны баримтаар 80--85 °C-д Arm цөмийн давтамж аажмаар буурна\strut
\end{minipage}\tabularnewline
\begin{minipage}[t]{0.30\columnwidth}\raggedright
microSD\strut
\end{minipage} & \begin{minipage}[t]{0.30\columnwidth}\raggedright
\textbf{A1 эсвэл A2}, 16 GB+\strut
\end{minipage} & \begin{minipage}[t]{0.30\columnwidth}\raggedright
Хямд карт дээр persistence, swap бичилт удаан\strut
\end{minipage}\tabularnewline
\begin{minipage}[t]{0.30\columnwidth}\raggedright
Сүлжээ\strut
\end{minipage} & \begin{minipage}[t]{0.30\columnwidth}\raggedright
\textbf{Кабель} (Wi-Fi биш)\strut
\end{minipage} & \begin{minipage}[t]{0.30\columnwidth}\raggedright
Pi 3B-ийн Wi-Fi нь зөвхөн 2.4 GHz, лабораторид бөглөрдөг\strut
\end{minipage}\tabularnewline
\bottomrule
\end{longtable}\endgroup

Pi 3B-ийн албан ёсны үзүүлэлт: BCM2837, 4 × Cortex-A53 (Armv8) @ 1.2 GHz, \textbf{1 GB} RAM, 100 Mb/s Ethernet, 4 × USB 2.0, 2.4 GHz 802.11n Wi-Fi, microSD.

\hypertarget{ux431.1-microsd-ux43aux430ux440ux442-ux431ux438ux447ux438ux445-raspberry-pi-imager}{%
\subsection{Б.1 microSD карт бичих (Raspberry Pi Imager)}\label{ux431.1-microsd-ux43aux430ux440ux442-ux431ux438ux447ux438ux445-raspberry-pi-imager}}

\begin{quote}
⚠ Бичихэд картын \textbf{бүх өгөгдөл устна}. Компьютерийн өөр диск эсвэл USB flash-ийг андуурч сонгохгүйн тулд \textbf{бусад USB төхөөрөмжөө салга}.
\end{quote}

\begin{enumerate}
\def\labelenumi{\arabic{enumi}.}
\item
  microSD картаа уншигчаар зөөврийн компьютерт залга.
  \textgreater{} ⚠ Windows \textbf{"You need to format the disk in drive X: before you can use it"} гэж асуувал \textbf{Cancel} дар. Format хийх ШААРДЛАГАГҮЙ.
\item
  \textbf{Raspberry Pi Imager}-ийг нээ (Windows зөвшөөрөл асуувал \textbf{Yes}).
\item
  \textbf{Device} → \textbf{Raspberry Pi 3} → \textbf{Next}.
\item
  \textbf{OS} → \textbf{Raspberry Pi OS (other)} → \textbf{Raspberry Pi OS Lite (64-bit)} → \textbf{Next}.

  \begin{quote}
  ⚠ \textbf{64-bit заавал.} Манай бүх image ба Python wheel (\texttt{ai-edge-litert} г.м.) arm64-д зориулагдсан. Docker Engine v29-өөс эхлэн Raspberry Pi OS 32-bit (armhf)-д шинэ багц гаргахаа больсон.

  ⚠ \textbf{Lite заавал.} Lite бол дэлгэцгүй, \emph{"command-line-only"} хувилбар --- ширээний орчин 1 GB-д хэт үнэтэй. Одоогийн хувилбар нь \textbf{Debian Trixie} дээр суурилсан (өмнөх нь Bookworm).
  \end{quote}
\item
  \textbf{Storage} → \textbf{microSD картаа} сонго. \textbf{Хэмжээг нь шалга} (16 GB карт бол \textasciitilde15 GB гэж харагдана) → \textbf{Next}.
\item
  \textbf{Customisation} (OS тохируулга) --- хуудас бүрт доорх утгыг оруулаад \textbf{Next}:

  \begingroup\scriptsize\begin{longtable}[]{@{}lll@{}}
  \toprule
  \begin{minipage}[b]{0.30\columnwidth}\raggedright
  Хуудас\strut
  \end{minipage} & \begin{minipage}[b]{0.30\columnwidth}\raggedright
  Талбар\strut
  \end{minipage} & \begin{minipage}[b]{0.30\columnwidth}\raggedright
  Утга\strut
  \end{minipage}\tabularnewline
  \midrule
  \endhead
  \begin{minipage}[t]{0.30\columnwidth}\raggedright
  \textbf{Hostname}\strut
  \end{minipage} & \begin{minipage}[t]{0.30\columnwidth}\raggedright
  Hostname\strut
  \end{minipage} & \begin{minipage}[t]{0.30\columnwidth}\raggedright
  \texttt{pi-team\textless{}NN\textgreater{}} (жишээ: \texttt{pi-team07})\strut
  \end{minipage}\tabularnewline
  \begin{minipage}[t]{0.30\columnwidth}\raggedright
  \textbf{Localisation}\strut
  \end{minipage} & \begin{minipage}[t]{0.30\columnwidth}\raggedright
  Capital city\strut
  \end{minipage} & \begin{minipage}[t]{0.30\columnwidth}\raggedright
  \textbf{Ulaanbaatar (Mongolia)}\strut
  \end{minipage}\tabularnewline
  \begin{minipage}[t]{0.30\columnwidth}\raggedright
  \strut
  \end{minipage} & \begin{minipage}[t]{0.30\columnwidth}\raggedright
  Time zone\strut
  \end{minipage} & \begin{minipage}[t]{0.30\columnwidth}\raggedright
  \texttt{Asia/Ulaanbaatar} (автоматаар бөглөгдөнө)\strut
  \end{minipage}\tabularnewline
  \begin{minipage}[t]{0.30\columnwidth}\raggedright
  \strut
  \end{minipage} & \begin{minipage}[t]{0.30\columnwidth}\raggedright
  Keyboard layout\strut
  \end{minipage} & \begin{minipage}[t]{0.30\columnwidth}\raggedright
  \texttt{us}\strut
  \end{minipage}\tabularnewline
  \begin{minipage}[t]{0.30\columnwidth}\raggedright
  \textbf{User}\strut
  \end{minipage} & \begin{minipage}[t]{0.30\columnwidth}\raggedright
  Username\strut
  \end{minipage} & \begin{minipage}[t]{0.30\columnwidth}\raggedright
  \textbf{\texttt{cnc302}} (өөр нэр бүү ашигла)\strut
  \end{minipage}\tabularnewline
  \begin{minipage}[t]{0.30\columnwidth}\raggedright
  \strut
  \end{minipage} & \begin{minipage}[t]{0.30\columnwidth}\raggedright
  Password\strut
  \end{minipage} & \begin{minipage}[t]{0.30\columnwidth}\raggedright
  Өөрийн нууц үг --- \textbf{тэмдэглэж ав} (\texttt{sudo}-д хэрэгтэй)\strut
  \end{minipage}\tabularnewline
  \begin{minipage}[t]{0.30\columnwidth}\raggedright
  \textbf{Wi-Fi}\strut
  \end{minipage} & \begin{minipage}[t]{0.30\columnwidth}\raggedright
  \strut
  \end{minipage} & \begin{minipage}[t]{0.30\columnwidth}\raggedright
  \textbf{Кабель} ашиглах бол алгас (\textbf{Skip}). Wi-Fi бол: SSID ба нууц үг (\textbf{2.4 GHz}, том/жижиг үсгийг яг зөв)\strut
  \end{minipage}\tabularnewline
  \begin{minipage}[t]{0.30\columnwidth}\raggedright
  \textbf{Remote access}\strut
  \end{minipage} & \begin{minipage}[t]{0.30\columnwidth}\raggedright
  Enable SSH\strut
  \end{minipage} & \begin{minipage}[t]{0.30\columnwidth}\raggedright
  \textbf{асаана}\strut
  \end{minipage}\tabularnewline
  \begin{minipage}[t]{0.30\columnwidth}\raggedright
  \strut
  \end{minipage} & \begin{minipage}[t]{0.30\columnwidth}\raggedright
  Authentication\strut
  \end{minipage} & \begin{minipage}[t]{0.30\columnwidth}\raggedright
  \textbf{Use public key authentication}\strut
  \end{minipage}\tabularnewline
  \begin{minipage}[t]{0.30\columnwidth}\raggedright
  \strut
  \end{minipage} & \begin{minipage}[t]{0.30\columnwidth}\raggedright
  Public key\strut
  \end{minipage} & \begin{minipage}[t]{0.30\columnwidth}\raggedright
  А.7, Алхам 2-т гарсан \textbf{бүтэн мөрийг} (\texttt{ssh-ed25519\ AAAA\ldots{}\ таны@имэйл.\allowbreak{}mn}) буулга\strut
  \end{minipage}\tabularnewline
  \begin{minipage}[t]{0.30\columnwidth}\raggedright
  \textbf{Raspberry Pi Connect}\strut
  \end{minipage} & \begin{minipage}[t]{0.30\columnwidth}\raggedright
  \strut
  \end{minipage} & \begin{minipage}[t]{0.30\columnwidth}\raggedright
  \textbf{Унтраалттай} үлдээ\strut
  \end{minipage}\tabularnewline
  \bottomrule
  \end{longtable}\endgroup

  \begin{quote}
  ⚠ Hostname \textbf{баг бүрт өөр} байх ёстой. Ижил нэртэй хоёр Pi нэг сүлжээнд байвал хооронд нь ялгах боломжгүй.

  ☞ Imager-ийн хувилбараас хамаарч хуудасны нэр бага зэрэг өөр байж болно --- утга нь адилхан.
  \end{quote}
\item
  \textbf{Write} → анхааруулга гарвал \textbf{I understand, erase and write} (эсвэл \textbf{Yes}) дар. Windows-ийн администраторын зөвшөөрөл асуувал \textbf{Yes}.
\item
  Бичих (\textbf{Writing}) ба баталгаажуулах (\textbf{Verifying}) явц дуусахыг хүлээ (5--15 мин). \textbf{Write complete!} гарвал картаа салга.
\end{enumerate}

\begin{quote}
Албан ёсны баримт: \href{https://www.raspberrypi.com/documentation/computers/getting-started.html}{Getting started --- Install using Imager} · \href{https://www.raspberrypi.com/documentation/computers/os.html}{Raspberry Pi OS} · \href{https://www.raspberrypi.com/news/a-new-raspberry-pi-imager/}{A new Raspberry Pi Imager (2.\allowbreak{}0)}
\end{quote}

\hypertarget{ux431.2-ux430ux43dux445-ux430ux441ux430ux430ux445-ssh-ux44dux44dux440-ux445ux43eux43bux431ux43eux433ux434ux43eux445-ux448ux438ux43dux44dux447ux43bux44dux445}{%
\subsection{Б.2 Анх асаах, SSH-ээр холбогдох, шинэчлэх}\label{ux431.2-ux430ux43dux445-ux430ux441ux430ux430ux445-ssh-ux44dux44dux440-ux445ux43eux43bux431ux43eux433ux434ux43eux445-ux448ux438ux43dux44dux447ux43bux44dux445}}

\textbf{Алхам 1 --- асаах.}

\begin{enumerate}
\def\labelenumi{\arabic{enumi}.}
\tightlist
\item
  microSD картаа Pi-д хий (контактаараа дээш биш, \textbf{доош} --- PCB тал руу).
\item
  Ethernet кабелиа зөөврийн компьютертэй \textbf{нэг router/switch}-д залга.
\item
  Тэжээлээ \textbf{хамгийн сүүлд} залга. Улаан гэрэл асаж, ногоон гэрэл анивчина.
\item
  \textbf{3--5 минут хүлээ.} Pi анх асахдаа файлын системээ өргөтгөж, өөрөө нэг-хоёр удаа restart хийнэ.
\end{enumerate}

\textbf{Алхам 2 --- Pi-г сүлжээнээс олох.}

\textbf{{[}ЗК · PowerShell{]}}

\begin{Shaded}
\begin{Highlighting}[]
\NormalTok{ping {-}4 pi{-}team\textless{}NN\textgreater{}.}\FunctionTok{local}
\end{Highlighting}
\end{Shaded}

✓ \texttt{Reply\ from\ 192.\allowbreak{}168.\allowbreak{}x.\allowbreak{}x:\allowbreak{}\ bytes=\allowbreak{}32\ time=\allowbreak{}\ldots{}} гарна. Тэр хаяг бол \texttt{\textless{}PI\_\allowbreak{}IP\textgreater{}} --- \textbf{тэмдэглэ}.

✗ \texttt{Ping\ request\ could\ not\ find\ host}: (а) дахиад 2 мин хүлээгээд оролд; (б) router-ийн удирдлагын хуудаснаас \textbf{connected devices / DHCP clients} жагсаалтад \texttt{pi-team\textless{}NN\textgreater{}}-ийг хайж IP-г ол; (в) Д хэсэг.

\textbf{Алхам 3 --- SSH-ээр холбогдох.}

\textbf{{[}ЗК · PowerShell{]}}

\begin{Shaded}
\begin{Highlighting}[]
\NormalTok{ssh pi}
\end{Highlighting}
\end{Shaded}

(\texttt{.local} ажиллахгүй бол: \texttt{ssh\ -i\ "\$env:USERPROFILE\textbackslash{}.ssh\textbackslash{}cnc302"\ cnc302@\textless{}PI\_IP\textgreater{}}.)

Анх удаа:

\begin{verbatim}
Are you sure you want to continue connecting (yes/no/[fingerprint])?
\end{verbatim}

→ \texttt{yes} гэж \textbf{бүтнээр} бичээд \texttt{Enter}.

✓ Нууц үг асуулгүйгээр prompt \texttt{cnc302@pi-team07:\textasciitilde{}\ \$} болж өөрчлөгдөнө. Одооноос \textbf{{[}Pi{]}} командыг \textbf{энэ цонхонд} бичнэ.

✗ \texttt{Permission\ denied\ (publickey)} → Д хэсэг.

\textbf{Алхам 4 --- системийг шинэчлэх ба хэрэгслүүд суулгах} (10--20 мин):

\textbf{{[}Pi{]}}

\begin{Shaded}
\begin{Highlighting}[]
\FunctionTok{sudo}\NormalTok{ apt update}
\FunctionTok{sudo}\NormalTok{ apt full{-}upgrade {-}y}
\FunctionTok{sudo}\NormalTok{ apt install {-}y git curl jq bc htop iotop stress{-}ng mosquitto{-}clients }\KeywordTok{\textbackslash{}}
                    \ExtensionTok{netcat{-}openbsd}\NormalTok{ python3{-}venv python3{-}dev}
\FunctionTok{sudo}\NormalTok{ reboot}
\end{Highlighting}
\end{Shaded}

\texttt{sudo} анх удаа Б.1-д тавьсан \textbf{Pi-гийн нууц үгийг} асууна. Сүүлийн мөрийн төгсгөлийн \texttt{\textbackslash{}} нь "команд дараагийн мөрөнд үргэлжилнэ" гэсэн утгатай --- хоёр мөрийг \textbf{хамт} хуул.

\texttt{reboot} хийхэд SSH тасарна (\texttt{Connection\ to\ \ldots{}\ closed}) --- энэ хэвийн. \textbf{1--2 мин хүлээгээд} \texttt{ssh\ pi}-ээр дахин холбогд.

\textbf{Алхам 5 --- үндсэн шалгалт.}

\textbf{{[}Pi{]}}

\begin{Shaded}
\begin{Highlighting}[]
\FunctionTok{uname}\NormalTok{ {-}m}
\ExtensionTok{dpkg}\NormalTok{ {-}{-}print{-}architecture}
\BuiltInTok{.} \ExtensionTok{/etc/os{-}release} \KeywordTok{\&\&} \BuiltInTok{echo} \StringTok{"}\VariableTok{$PRETTY\_NAME}\StringTok{"}
\FunctionTok{free}\NormalTok{ {-}h}
\ExtensionTok{vcgencmd}\NormalTok{ get\_throttled}
\end{Highlighting}
\end{Shaded}

✓ \textbf{Хүлээгдэх үр дүн:}

\begingroup\scriptsize\begin{longtable}[]{@{}lll@{}}
\toprule
Команд & Байх ёстой & Өөр бол\tabularnewline
\midrule
\endhead
\texttt{uname\ -m} & \texttt{aarch64} & \texttt{armv7l} = 32-bit OS → Б.1\tabularnewline
\texttt{dpkg\ -\/\allowbreak{}-print-architecture} & \texttt{arm64} & дээрхтэй адил\tabularnewline
\texttt{PRETTY\_NAME} & \texttt{\ldots{}\ 13\ (trixie)} эсвэл \texttt{\ldots{}\ 12\ (bookworm)} & ---\tabularnewline
\texttt{free\ -h} → Mem total & ≈ \texttt{900Mi} & ---\tabularnewline
\texttt{get\_throttled} & \texttt{throttled=0x0} & Тэжээл сул → Б.0, Д\tabularnewline
\bottomrule
\end{longtable}\endgroup

\hypertarget{ux431.3-gpu-ux433ux438ux439ux43d-ux441ux430ux43dux430ux445-ux43eux439-gpu_mem-ux445ux44dux43cux436ux438ux436-ux448ux438ux439ux434ux43dux44d}{%
\subsection{\texorpdfstring{Б.3 GPU-гийн санах ой (\texttt{gpu\_mem}) --- хэмжиж шийднэ}{Б.3 GPU-гийн санах ой (gpu\_mem) --- хэмжиж шийднэ}}\label{ux431.3-gpu-ux433ux438ux439ux43d-ux441ux430ux43dux430ux445-ux43eux439-gpu_mem-ux445ux44dux43cux436ux438ux436-ux448ux438ux439ux434ux43dux44d}}

Raspberry Pi-гийн баримтаар 1 GB-тай загварт \texttt{gpu\_mem}-ийн \textbf{анхдагч утга 76} MiB. Гэхдээ баримт \texttt{gpu\_mem}-ийг \textbf{legacy} тохиргоонд ангилж, \emph{Raspberry Pi OS Bookworm болон түүнээс хойшхи хувилбарт ажиллахгүй, албан ёсоор дэмжигдэхгүй} гэж тэмдэглэсэн. Тиймээс «N MiB хэмнэнэ» гэж амлахгүй --- \textbf{өмнө ба дараа нь хэмжинэ}.

\textbf{Алхам 1 --- өмнөх утгыг хэмжих.}

\textbf{{[}Pi{]}}

\begin{Shaded}
\begin{Highlighting}[]
\FunctionTok{grep}\NormalTok{ MemTotal /proc/meminfo}
\ExtensionTok{vcgencmd}\NormalTok{ get\_mem gpu}
\end{Highlighting}
\end{Shaded}

Хоёр утгыг (жишээ: \texttt{MemTotal:\allowbreak{}\ 925\ldots{}kB}, \texttt{gpu=76M}) \textbf{тэмдэглэ} --- (1).

\textbf{Алхам 2 --- тохиргоо нэмж, дахин ачаалах.}

\textbf{{[}Pi{]}}

\begin{Shaded}
\begin{Highlighting}[]
\BuiltInTok{echo} \StringTok{\textquotesingle{}gpu\_mem=16\textquotesingle{}} \KeywordTok{|} \FunctionTok{sudo}\NormalTok{ tee {-}a /boot/firmware/config.txt}
\FunctionTok{sudo}\NormalTok{ reboot}
\end{Highlighting}
\end{Shaded}

\textbf{Алхам 3 --- дахин холбогдоод} (\texttt{ssh\ pi}) \textbf{хэмжих.}

\textbf{{[}Pi{]}}

\begin{Shaded}
\begin{Highlighting}[]
\FunctionTok{grep}\NormalTok{ MemTotal /proc/meminfo}
\ExtensionTok{vcgencmd}\NormalTok{ get\_mem gpu}
\end{Highlighting}
\end{Shaded}

Утгыг \textbf{тэмдэглэ} --- (2). (1) ба (2)-ын зөрүүг Лаб 1-ийн хүснэгтэд бич. Зөрүүгүй бол энэ мөр таны OS дээр нөлөөгүй гэсэн үг --- буруу биш, \textbf{хэмжилт}.

\begin{quote}
Албан ёсны баримт: \href{https://www.raspberrypi.com/documentation/computers/legacy_config_txt.html\#gpu_mem}{Legacy config.\allowbreak{}txt --- gpu\_\allowbreak{}mem}
\end{quote}

\hypertarget{ux431.4-docker-engine-ux441ux443ux443ux43bux433ux430ux445}{%
\subsection{Б.4 Docker Engine суулгах}\label{ux431.4-docker-engine-ux441ux443ux443ux43bux433ux430ux445}}

Pi дээр Docker \textbf{Desktop} биш, \textbf{Docker Engine} суулгана. 64-бит Raspberry Pi OS нь Debian-д суурилсан тул Docker-ийн \textbf{Debian} зааврыг (\texttt{apt} repository) дагана. \texttt{get.docker.com} convenience script-ийг Docker өөрөө \emph{зөвхөн туршилт ба хөгжүүлэлтийн орчинд} зөвлөдөг тул бид ашиглахгүй.

\textbf{Алхам 1 --- зөрчилдөх албан бус багцыг устгах.} Шинэ систем дээр эдгээр байхгүй тул \texttt{no\ packages\ found} эсвэл \texttt{0\ to\ remove} гарна --- \textbf{энэ хэвийн}.

\textbf{{[}Pi{]}}

\begin{Shaded}
\begin{Highlighting}[]
\FunctionTok{sudo}\NormalTok{ apt remove }\VariableTok{$(}\ExtensionTok{dpkg}\NormalTok{ {-}{-}get{-}selections docker.io docker{-}compose docker{-}doc docker{-}buildx podman{-}docker containerd runc }\KeywordTok{|} \FunctionTok{cut}\NormalTok{ {-}f1}\VariableTok{)}
\end{Highlighting}
\end{Shaded}

\textbf{Алхам 2 --- Docker-ийн GPG түлхүүр нэмэх} (\texttt{apt} Docker-ийн багцыг жинхэнэ эсэхийг үүгээр шалгана):

\textbf{{[}Pi{]}}

\begin{Shaded}
\begin{Highlighting}[]
\FunctionTok{sudo}\NormalTok{ apt update}
\FunctionTok{sudo}\NormalTok{ apt install {-}y ca{-}certificates curl}
\FunctionTok{sudo}\NormalTok{ install {-}m 0755 {-}d /etc/apt/keyrings}
\FunctionTok{sudo}\NormalTok{ curl {-}fsSL https://download.docker.com/linux/debian/gpg {-}o /etc/apt/keyrings/docker.asc}
\FunctionTok{sudo}\NormalTok{ chmod a+r /etc/apt/keyrings/docker.asc}
\end{Highlighting}
\end{Shaded}

✓ \textbf{Шалгах:} \texttt{ls\ -l\ /\allowbreak{}etc/\allowbreak{}apt/\allowbreak{}keyrings/\allowbreak{}docker.\allowbreak{}asc} → файл байна (хэмжээ \textasciitilde3.8 KB).

\textbf{Алхам 3 --- Docker-ийн repository нэмэх.} Доорх блокийг \texttt{sudo\ tee}-ээс \texttt{EOF} хүртэл \textbf{бүтнээр нь НЭГ ДОР} хуулж буулга:

\textbf{{[}Pi{]}}

\begin{Shaded}
\begin{Highlighting}[]
\FunctionTok{sudo}\NormalTok{ tee /etc/apt/sources.list.d/docker.sources }\OperatorTok{\textless{}\textless{}EOF}
\NormalTok{Types: deb}
\NormalTok{URIs: https://download.docker.com/linux/debian}
\NormalTok{Suites: }\VariableTok{$(}\BuiltInTok{.} \ExtensionTok{/etc/os{-}release} \KeywordTok{\&\&} \BuiltInTok{echo} \StringTok{"}\VariableTok{$VERSION\_CODENAME}\StringTok{"}\VariableTok{)}
\NormalTok{Components: stable}
\NormalTok{Architectures: }\VariableTok{$(}\ExtensionTok{dpkg}\NormalTok{ {-}{-}print{-}architecture}\VariableTok{)}
\NormalTok{Signed{-}By: /etc/apt/keyrings/docker.asc}
\NormalTok{EOF}
\end{Highlighting}
\end{Shaded}

✓ \textbf{Шалгах:} терминал файлын агуулгыг хэвлэнэ. Тэнд \textbf{\texttt{Suites:\ trixie}} (эсвэл \texttt{bookworm}) ба \textbf{\texttt{Architectures:\allowbreak{}\ arm64}} байх ёстой. Хоосон бол Алхам 3-ыг дахин хий.

\textbf{Алхам 4 --- суулгах.}

\textbf{{[}Pi{]}}

\begin{Shaded}
\begin{Highlighting}[]
\FunctionTok{sudo}\NormalTok{ apt update}
\FunctionTok{sudo}\NormalTok{ apt install {-}y docker{-}ce docker{-}ce{-}cli containerd.io docker{-}buildx{-}plugin docker{-}compose{-}plugin}
\end{Highlighting}
\end{Shaded}

✓ \texttt{sudo\ apt\ update}-ийн гаралтад \texttt{https:\allowbreak{}/\allowbreak{}/\allowbreak{}download.\allowbreak{}docker.\allowbreak{}com/\allowbreak{}linux/\allowbreak{}debian\ trixie\ InRelease} мөр харагдана. Суулгалт 3--5 мин.

\textbf{Алхам 5 --- \texttt{sudo}-гүй ажиллуулах эрх өгөх.}

\textbf{{[}Pi{]}}

\begin{Shaded}
\begin{Highlighting}[]
\FunctionTok{sudo}\NormalTok{ usermod {-}aG docker }\VariableTok{$USER}
\BuiltInTok{exit}
\end{Highlighting}
\end{Shaded}

\begin{quote}
⚠ \texttt{exit} хийж SSH-ээс \textbf{заавал гараад}, \texttt{ssh\ pi}-ээр \textbf{дахин холбогд}. Бүлгийн эрх зөвхөн шинэ session-д хүчинтэй болно.

⚠ \texttt{docker} бүлэгт орсон хэрэглэгч root-той дүйцэх эрхтэй болдог (Docker-ийн post-install баримт) --- Pi-гийн \texttt{cnc302} хэрэглэгчийн нууц үгийг хамгаал.
\end{quote}

\textbf{Алхам 6 --- шалгах.}

\textbf{{[}Pi{]}}

\begin{Shaded}
\begin{Highlighting}[]
\ExtensionTok{docker}\NormalTok{ {-}{-}version}
\ExtensionTok{docker}\NormalTok{ compose version}
\ExtensionTok{docker}\NormalTok{ run {-}{-}rm hello{-}world}
\ExtensionTok{systemctl}\NormalTok{ is{-}enabled docker}
\end{Highlighting}
\end{Shaded}

✓ \textbf{Хүлээгдэх үр дүн:}

\begin{itemize}
\tightlist
\item
  \texttt{Docker\ version\ 2x.\allowbreak{}x.\allowbreak{}x} ба \texttt{Docker\ Compose\ version\ v2.\allowbreak{}x.\allowbreak{}x} (\textbf{v2.20+}),
\item
  \texttt{Hello\ from\ Docker!},
\item
  \texttt{enabled} (Pi асах бүрт Docker автоматаар асна).
\end{itemize}

✗ \texttt{permission\ denied\ while\ trying\ to\ connect\ to\ the\ Docker\ daemon\ socket} → Алхам 5-ын дараа SSH-ээс гараагүй байна.

\textbf{Алхам 7 --- Docker-ийн лог microSD-г элээхээс сэргийлэх.} Блокийг \textbf{бүтнээр нь} буулга (\texttt{log-opts}-ийн утгууд \textbf{мөр} --- хашилтанд --- байх ёстой):

\textbf{{[}Pi{]}}

\begin{Shaded}
\begin{Highlighting}[]
\FunctionTok{sudo}\NormalTok{ tee /etc/docker/daemon.json }\OperatorTok{\textgreater{}}\NormalTok{ /dev/null }\OperatorTok{\textless{}\textless{}\textquotesingle{}EOF\textquotesingle{}}
\NormalTok{\{}
\NormalTok{  "log{-}driver": "json{-}file",}
\NormalTok{  "log{-}opts": \{ "max{-}size": "5m", "max{-}file": "2" \}}
\NormalTok{\}}
\OperatorTok{EOF}
\FunctionTok{sudo}\NormalTok{ systemctl restart docker}
\end{Highlighting}
\end{Shaded}

✓ \textbf{Шалгах:} \texttt{docker\ info\ -\/-format\ \textquotesingle{}\{\{.LoggingDriver\}\}\textquotesingle{}} → \texttt{json-file}. Энэ тохиргоо зөвхөн \textbf{шинээр үүсэх} контейнерт үйлчилнэ.

✗ \texttt{systemctl\ restart\ docker} алдаа гаргавал JSON-д алдаа байна: \texttt{sudo\ nano\ /\allowbreak{}etc/\allowbreak{}docker/\allowbreak{}daemon.\allowbreak{}json}-оор дээрх блоктой тулгаж засна.

\begin{quote}
Албан ёсны баримт: \href{https://docs.docker.com/engine/install/debian/}{Install Docker Engine on Debian} · \href{https://docs.docker.com/engine/install/linux-postinstall/}{Linux post-installation steps} · \href{https://docs.docker.com/engine/logging/drivers/json-file/}{JSON File logging driver}
\end{quote}

\hypertarget{ux431.5-ux441ux43eux43bux438ux445-ux441ux430ux43dux430ux445-ux43eux439-swap-ux431ux430-cgroup-ux437ux430ux430ux432ux430ux43b}{%
\subsection{Б.5 Солих санах ой (swap) ба cgroup (ЗААВАЛ)}\label{ux431.5-ux441ux43eux43bux438ux445-ux441ux430ux43dux430ux445-ux43eux439-swap-ux431ux430-cgroup-ux437ux430ux430ux432ux430ux43b}}

\hypertarget{swap}{%
\subsubsection{Swap}\label{swap}}

1 GB RAM дээр swap нь Лаб 4--8-д OOM-оос хамгаалах \textbf{аюулгүйн сүлжээ}. Курсын бүх заавар \textbf{2 GiB} swap гэж тооцсон. Raspberry Pi OS-ийн хувилбараас хамаарч swap-ыг удирдах хэрэгсэл өөр байдаг тул \textbf{эхлээд аль нь болохыг} тодорхойлно.

\textbf{Алхам 1 --- одоогийн байдал.}

\textbf{{[}Pi{]}}

\begin{Shaded}
\begin{Highlighting}[]
\ExtensionTok{swapon}\NormalTok{ {-}{-}show}\KeywordTok{;} \FunctionTok{free}\NormalTok{ {-}h}
\BuiltInTok{test}\NormalTok{ {-}e /etc/rpi/swap.conf }\KeywordTok{\&\&} \BuiltInTok{echo} \StringTok{"rpi{-}swap"} \KeywordTok{||} \BuiltInTok{echo} \StringTok{"rpi{-}swap алга"}
\ExtensionTok{systemctl}\NormalTok{ is{-}enabled dphys{-}swapfile }\OperatorTok{2\textgreater{}}\NormalTok{/dev/null}
\end{Highlighting}
\end{Shaded}

\begingroup\scriptsize\begin{longtable}[]{@{}ll@{}}
\toprule
Хоёр дахь мөрийн гаралт & Аль замаар явах\tabularnewline
\midrule
\endhead
\texttt{rpi-swap} & \textbf{Алхам 2а} (шинэ Trixie image --- ихэнх тохиолдол)\tabularnewline
\texttt{rpi-swap\ алга} ба дараа нь \texttt{enabled} & \textbf{Алхам 2б} (хуучин Bookworm image)\tabularnewline
\bottomrule
\end{longtable}\endgroup

\textbf{Алхам 2а --- \texttt{rpi-swap} байгаа бол.} Raspberry Pi-гийн \texttt{rpi-swap} багц \texttt{dphys-swapfile}-ийг орлох зорилготой; анхдагч \texttt{Mechanism=auto} нь одоогоор \textbf{zram+file} (шахсан RAM swap + ховор бичигддэг файл), zram-ын хэмжээ анхдагчаар RAM-тай тэнцүү (2048 MiB-ээр хязгаарлагдана) --- Pi 3B дээр \textasciitilde0.9 GiB. Бид \texttt{swap.conf(5)}-ын \textbf{албан ёсны "Example 1"}-ийг дагаж 2 GiB-ийн уламжлалт swap файл тавина. Блокийг \textbf{бүтнээр нь} буулга:

\textbf{{[}Pi{]}}

\begin{Shaded}
\begin{Highlighting}[]
\FunctionTok{sudo}\NormalTok{ mkdir {-}p /etc/rpi/swap.conf.d/}
\FunctionTok{sudo}\NormalTok{ tee /etc/rpi/swap.conf.d/80{-}use{-}swapfile.conf }\OperatorTok{\textgreater{}}\NormalTok{ /dev/null }\OperatorTok{\textless{}\textless{}EOF}
\NormalTok{[Main]}
\NormalTok{Mechanism=swapfile}

\NormalTok{[File]}
\NormalTok{FixedSizeMiB=2048}
\OperatorTok{EOF}
\FunctionTok{sudo}\NormalTok{ reboot}
\end{Highlighting}
\end{Shaded}

swap-ын тохиргоо \textbf{зөвхөн дахин ачаалахад} хэрэгжинэ. 1--2 мин хүлээгээд \texttt{ssh\ pi}.

\begin{quote}
Анхдагч zram+file-ийг үлдээх нь microSD-г бага элээнэ (албан ёсны зорилго нь тэр). Энэ тохиолдолд \texttt{make\ check}-ийн swap мөр ✗ гарах бөгөөд Лаб 1--8-ын "2.0Gi" гэсэн хүлээлтийн оронд өөрийн утгаа тайландаа бич.
\end{quote}

\textbf{Алхам 2б --- \texttt{rpi-swap} байхгүй, \texttt{dphys-swapfile} ажиллаж байгаа бол.} Энэ багцын тохиргоо Raspberry Pi-гийн баримтад ороогүй тул доорх нь уламжлалт арга:

\textbf{{[}Pi{]}}

\begin{Shaded}
\begin{Highlighting}[]
\FunctionTok{sudo}\NormalTok{ dphys{-}swapfile swapoff}
\FunctionTok{sudo}\NormalTok{ sed {-}i }\StringTok{\textquotesingle{}s/\^{}\#\textbackslash{}?CONF\_SWAPSIZE=.*/CONF\_SWAPSIZE=2048/\textquotesingle{}}\NormalTok{ /etc/dphys{-}swapfile}
\FunctionTok{sudo}\NormalTok{ dphys{-}swapfile setup}
\FunctionTok{sudo}\NormalTok{ dphys{-}swapfile swapon}
\end{Highlighting}
\end{Shaded}

\textbf{Алхам 3 --- шалгах} (аль ч замаар явсан):

\textbf{{[}Pi{]}}

\begin{Shaded}
\begin{Highlighting}[]
\FunctionTok{free}\NormalTok{ {-}h}
\end{Highlighting}
\end{Shaded}

✓ \texttt{Swap:} мөрийн \textbf{total} баганад \texttt{2.0Gi}.

\begin{quote}
\textbf{Санамж:} swap файл microSD дээр байгаа тул \textbf{маш удаан}. Энэ бол гүйцэтгэлийн шийдэл биш. Лаб 4-д swap ажиллаж эхэлбэл (\texttt{vmstat\ 1}-ийн \texttt{si/so} багана 0-ээс их) та аль хэдийн хязгаарт хүрсэн --- тэр нь хэмжилт өөрөө юм.

Албан ёсны эх: \href{https://github.com/raspberrypi/rpi-swap}{raspberrypi/\allowbreak{}rpi-swap --- README, swap.\allowbreak{}conf(5)}
\end{quote}

\hypertarget{cgroup}{%
\subsubsection{cgroup}\label{cgroup}}

K3s-ийн баримтаар стандарт Raspberry Pi OS cgroup-ийн санах ойн хянагчийг идэвхжүүлээгүй эхэлдэг; Лаб 8-ын K3s agent, \texttt{docker\ stats}-ын санах ойн тоо ба ирмэгийн агентын \texttt{MemoryMax=200M} хязгаарт хэрэгтэй. Параметрийг \texttt{/\allowbreak{}boot/\allowbreak{}firmware/\allowbreak{}cmdline.\allowbreak{}txt}-ийн \textbf{цорын ганц мөрийн төгсгөлд} залгана (Debian 11 ба түүнээс өмнөх image-д зам нь \texttt{/\allowbreak{}boot/\allowbreak{}cmdline.\allowbreak{}txt}).

\begin{quote}
⚠ Raspberry Pi-гийн баримт: \texttt{cmdline.txt}-ийн бүх параметр \textbf{нэг мөрөнд} байх ёстой --- шинэ мөр оруулбал түүнээс хойшхыг цөм үл тоомсорлоно. Тиймээс файлыг гараар \textbf{бүү} зас, доорх командыг ашигла.
\end{quote}

\textbf{Алхам 1 --- параметр нэмэх ба дахин ачаалах.}

\textbf{{[}Pi{]}}

\begin{Shaded}
\begin{Highlighting}[]
\FunctionTok{grep}\NormalTok{ {-}q }\StringTok{\textquotesingle{}cgroup\_memory=1\textquotesingle{}}\NormalTok{ /boot/firmware/cmdline.txt }\KeywordTok{||} \KeywordTok{\textbackslash{}}
  \FunctionTok{sudo}\NormalTok{ sed {-}i }\StringTok{\textquotesingle{}1 s/$/ cgroup\_memory=1 cgroup\_enable=memory/\textquotesingle{}}\NormalTok{ /boot/firmware/cmdline.txt}
\FunctionTok{cat}\NormalTok{ /boot/firmware/cmdline.txt}
\end{Highlighting}
\end{Shaded}

✓ Сүүлийн команд \textbf{НЭГ мөр} хэвлэх бөгөөд төгсгөлд нь \texttt{cgroup\_\allowbreak{}memory=\allowbreak{}1\ cgroup\_\allowbreak{}enable=\allowbreak{}memory} байна. Тэгвэл:

\textbf{{[}Pi{]}}

\begin{Shaded}
\begin{Highlighting}[]
\FunctionTok{sudo}\NormalTok{ reboot}
\end{Highlighting}
\end{Shaded}

\textbf{Алхам 2 --- дахин холбогдоод шалгах.}

\textbf{{[}Pi{]}}

\begin{Shaded}
\begin{Highlighting}[]
\FunctionTok{grep}\NormalTok{ {-}w memory /sys/fs/cgroup/cgroup.controllers}
\FunctionTok{free}\NormalTok{ {-}m }\KeywordTok{|} \FunctionTok{head}\NormalTok{ {-}2}
\end{Highlighting}
\end{Shaded}

✓ Эхний команд \texttt{memory} үгийг агуулсан мөр хэвлэнэ. \texttt{MemTotal}-ийг Лаб 1-д бичнэ.

\begin{quote}
Албан ёсны баримт: \href{https://docs.k3s.io/installation/requirements}{K3s Requirements --- Raspberry Pi (cgroups)} · \href{https://www.raspberrypi.com/documentation/computers/configuration.html\#configure-the-kernel-command-line}{Kernel command line (cmdline.\allowbreak{}txt)}
\end{quote}

\hypertarget{ux431.6-ux446ux430ux433ux438ux439ux43d-ux441ux438ux43dux445ux440ux43eux43dux447ux43bux43eux43b-ux437ux430ux430ux432ux430ux43b}{%
\subsection{Б.6 Цагийн синхрончлол (ЗААВАЛ)}\label{ux431.6-ux446ux430ux433ux438ux439ux43d-ux441ux438ux43dux445ux440ux43eux43dux447ux43bux43eux43b-ux437ux430ux430ux432ux430ux43b}}

Лаб 2-ын TLS гэрчилгээ, Лаб 3 ба Лаб 5-ын саатлын хэмжилт буруу цагтай үед утгагүй болдог. Pi 3B-д бодит цагийн цаг (RTC) \textbf{байхгүй} --- асаах бүрд NTP-ээс цагаа авдаг.

\textbf{{[}Pi{]}}

\begin{Shaded}
\begin{Highlighting}[]
\FunctionTok{sudo}\NormalTok{ timedatectl set{-}timezone Asia/Ulaanbaatar}
\FunctionTok{sudo}\NormalTok{ timedatectl set{-}ntp true}
\ExtensionTok{timedatectl}\NormalTok{ status}
\end{Highlighting}
\end{Shaded}

✓ \texttt{Time\ zone:\allowbreak{}\ Asia/\allowbreak{}Ulaanbaatar\ (+08,\ +0800)} ба \texttt{System\ clock\ synchronized:\allowbreak{}\ yes}.

✗ \texttt{synchronized:\allowbreak{}\ no} хэвээр бол 1--2 мин хүлээгээд дахин шалга; тэгээд ч болохгүй бол Pi интернэтэд гарч чадахгүй байна (\texttt{ping\ -c3\ 8.\allowbreak{}8.\allowbreak{}8.\allowbreak{}8}).

\hypertarget{ux431.7-ux441ux430ux43dux433ux430ux430-pi-ux434ux44dux44dux440-ux431ux430ux439ux440ux43bux443ux443ux43bux430ux445}{%
\subsection{Б.7 Сангаа Pi дээр байрлуулах}\label{ux431.7-ux441ux430ux43dux433ux430ux430-pi-ux434ux44dux44dux440-ux431ux430ux439ux440ux43bux443ux443ux43bux430ux445}}

\begin{quote}
⚠ Pi дээрх сангийн зам \textbf{заавал \texttt{\textasciitilde{}/cnc302}} (\texttt{/\allowbreak{}home/\allowbreak{}cnc302/\allowbreak{}cnc302}) байна --- \texttt{edge/\allowbreak{}agent/\allowbreak{}cnc302-edge-agent.\allowbreak{}service} ба лабын заавар яг энэ замыг хүлээнэ.
\end{quote}

\textbf{Алхам 1 --- clone хийх.}

\textbf{{[}Pi{]}}

\begin{Shaded}
\begin{Highlighting}[]
\BuiltInTok{cd}\NormalTok{ \textasciitilde{}}
\FunctionTok{git}\NormalTok{ clone }\OperatorTok{\textless{}}\NormalTok{REPO\_URL}\OperatorTok{\textgreater{}}\NormalTok{ cnc302}
\FunctionTok{ls}\NormalTok{ \textasciitilde{}/cnc302}
\end{Highlighting}
\end{Shaded}

\begin{quote}
☞ Багийн сан \textbf{private} бол GitHub HTTPS-ээр нэвтрэхэд нууц үгийн оронд \textbf{Personal Access Token} шаардана: \texttt{Username} → GitHub-ын нэр, \texttt{Password} → token (GitHub → Settings → Developer settings → Personal access tokens). Токенийг Pi дээр файлд бүү хадгал.
\end{quote}

\textbf{Алхам 2 --- ирмэгийн тохиргоо.}

\textbf{{[}Pi{]}}

\begin{Shaded}
\begin{Highlighting}[]
\BuiltInTok{cd}\NormalTok{ \textasciitilde{}/cnc302/edge}
\FunctionTok{cp}\NormalTok{ .env.example .env}
\FunctionTok{nano}\NormalTok{ .env}
\end{Highlighting}
\end{Shaded}

Дараах \textbf{хоёр мөрийг} соль (бусдыг хэвээр нь үлдээ):

\begingroup\scriptsize\begin{longtable}[]{@{}lll@{}}
\toprule
Мөр & Утга & Жишээ\tabularnewline
\midrule
\endhead
\texttt{CLOUD\_HOST=} & А.5-д олсон \texttt{\textless{}LAPTOP\_\allowbreak{}IP\textgreater{}} & \texttt{CLOUD\_\allowbreak{}HOST=\allowbreak{}192.\allowbreak{}168.\allowbreak{}1.\allowbreak{}100}\tabularnewline
\texttt{DEVICE\_ID=} & \texttt{pi3b-team\textless{}NN\textgreater{}} & \texttt{DEVICE\_\allowbreak{}ID=\allowbreak{}pi3b-team07}\tabularnewline
\bottomrule
\end{longtable}\endgroup

Хадгалах: \texttt{Ctrl\ +\ O} → \texttt{Enter} → \texttt{Ctrl\ +\ X}.

\begin{quote}
⚠ \texttt{.env}-д тайлбарыг \textbf{тусдаа мөрөнд} бич (\texttt{KEY=\allowbreak{}утга\ \ \#\ тайлбар} БИШ), утгыг хашилтгүй, хоосон зайгүй бич: systemd-ийн \texttt{EnvironmentFile=} мөрийн дундах \texttt{\#}-ийг утгын нэг хэсэг гэж уншдаг.
\end{quote}

\textbf{Алхам 3 --- гүүрний тохиргоо үүсгэх.}

\textbf{{[}Pi{]}}

\begin{Shaded}
\begin{Highlighting}[]
\FunctionTok{make}\NormalTok{ bridge}
\end{Highlighting}
\end{Shaded}

✓ \texttt{✓\ mosquitto/\allowbreak{}conf.\allowbreak{}d/\allowbreak{}bridge.\allowbreak{}conf\ үүслээ\ \ →\ \ 192.\allowbreak{}168.\allowbreak{}1.\allowbreak{}100:\allowbreak{}1883} (таны IP) гарна.

✗ \texttt{⚠\ CLOUD\_\allowbreak{}HOST\ хэвээрээ\ байна} → Алхам 2-т IP-г сольж хадгалаагүй байна.

\hypertarget{ux431.8-ux438ux440ux43cux44dux433ux438ux439ux43d-image-ux431ux430-python-ux43eux440ux447ux438ux43d}{%
\subsection{Б.8 Ирмэгийн image ба Python орчин}\label{ux431.8-ux438ux440ux43cux44dux433ux438ux439ux43d-image-ux431ux430-python-ux43eux440ux447ux438ux43d}}

\textbf{{[}Pi{]}}

\begin{Shaded}
\begin{Highlighting}[]
\BuiltInTok{cd}\NormalTok{ \textasciitilde{}/cnc302/edge}
\ExtensionTok{docker}\NormalTok{ compose pull}
\ExtensionTok{python3}\NormalTok{ {-}m venv .venv}
\ExtensionTok{.venv/bin/pip}\NormalTok{ install {-}{-}upgrade pip}
\ExtensionTok{.venv/bin/pip}\NormalTok{ install {-}r agent/requirements.txt}
\end{Highlighting}
\end{Shaded}

\texttt{pip\ install} Pi 3B дээр \textbf{5--15 мин} үргэлжилж магадгүй --- дэлгэц хэсэг хугацаанд өөрчлөгдөхгүй байсан ч \textbf{бүү тасал}.

✓ \textbf{Шалгах:}

\textbf{{[}Pi{]}}

\begin{Shaded}
\begin{Highlighting}[]
\ExtensionTok{.venv/bin/python}\NormalTok{ {-}c }\StringTok{"from ai\_edge\_litert.interpreter import Interpreter; print(\textquotesingle{}LiteRT OK\textquotesingle{})"}
\end{Highlighting}
\end{Shaded}

\texttt{LiteRT\ OK} гарна.

\begin{quote}
\texttt{ai-edge-litert} (LiteRT --- TensorFlow Lite-ийн шинэ нэр) нь PyPI дээр Python 3.11 (Bookworm) ба 3.13 (Trixie)-д зориулсан \texttt{manylinux\_\allowbreak{}2\_\allowbreak{}27\_\allowbreak{}aarch64} wheel-тэй; \texttt{numpy} ч мөн aarch64 wheel-тэй тул Pi дээр эх кодоос барихгүй. \texttt{No\ matching\ distribution} гарвал \texttt{uname\ -m} → \texttt{aarch64} эсэхийг шалга. Хуучин \texttt{tflite-runtime} (сүүлийн хувилбар 2.14, 2023) нь зөвхөн нөөц хувилбар.

Албан ёсны баримт: \href{https://ai.google.dev/edge/litert/migration}{LiteRT --- Migrate from tflite-runtime} · \href{https://pypi.org/project/ai-edge-litert/}{pypi: ai-edge-litert}
\end{quote}

\hypertarget{ux432.-ux437ux43api-ux445ux43eux451ux440-ux43cux430ux448ux438ux43dux44bux433-ux445ux43eux43bux431ux43eux445}{%
\section{В. {[}ЗК{]}{[}Pi{]} Хоёр машиныг холбох}\label{ux432.-ux437ux43api-ux445ux43eux451ux440-ux43cux430ux448ux438ux43dux44bux433-ux445ux43eux43bux431ux43eux445}}

Энэ бол бүх лабын суурь: Pi-гаас илгээсэн MQTT мессеж зөөврийн компьютер дээрх EMQX-д хүрэх ёстой. \textbf{Гурван терминал} зэрэг нээлттэй байна:

\begingroup\scriptsize\begin{longtable}[]{@{}lll@{}}
\toprule
Цонх & Терминал & Үүрэг\tabularnewline
\midrule
\endhead
1 & {[}Pi{]} & Гүүр ба агент\tabularnewline
2 & {[}Pi{]} (хоёр дахь \texttt{ssh\ pi}) & Гүүрний төлөв\tabularnewline
3 & {[}ЗК · Ubuntu{]} & Үүлэн дэх мессежийг сонсох\tabularnewline
\bottomrule
\end{longtable}\endgroup

\hypertarget{ux432.1-ux433ux4afux4afux440ux438ux439ux433-ux430ux441ux430ux430ux445}{%
\subsection{В.1 Гүүрийг асаах}\label{ux432.1-ux433ux4afux4afux440ux438ux439ux433-ux430ux441ux430ux430ux445}}

Эхлээд зөөврийн компьютер дээр стек ажиллаж байгааг шалга:

\textbf{{[}ЗК · Ubuntu{]}}

\begin{Shaded}
\begin{Highlighting}[]
\BuiltInTok{cd}\NormalTok{ \textasciitilde{}/cnc302/stack }\KeywordTok{\&\&} \FunctionTok{make}\NormalTok{ up }\KeywordTok{\&\&} \FunctionTok{make}\NormalTok{ health}
\end{Highlighting}
\end{Shaded}

Дараа нь:

\textbf{{[}Pi{]}} (цонх 1)

\begin{Shaded}
\begin{Highlighting}[]
\BuiltInTok{cd}\NormalTok{ \textasciitilde{}/cnc302/edge}
\FunctionTok{make}\NormalTok{ up}
\end{Highlighting}
\end{Shaded}

✓ \texttt{docker\ compose\ ps}-ын жагсаалтад \texttt{mosquitto} → \texttt{Up}.

\textbf{{[}Pi{]}} (цонх 2 --- PowerShell-ээс дахин \texttt{ssh\ pi})

\begin{Shaded}
\begin{Highlighting}[]
\BuiltInTok{cd}\NormalTok{ \textasciitilde{}/cnc302/edge}
\FunctionTok{make}\NormalTok{ link}
\end{Highlighting}
\end{Shaded}

✓ \texttt{cnc302/\allowbreak{}shutis/\allowbreak{}\ldots{}/\allowbreak{}bridge/\allowbreak{}state\ 1} гарвал \textbf{гүүр холбогдсон}. \texttt{0} бол тасарсан → В.2-ын хүснэгт. \texttt{Ctrl\ +\ C}-ээр зогсооно.

\hypertarget{ux432.2-ux445ux43eux451ux440-ux442ux430ux43bux430ux430ux441-ux448ux430ux43bux433ux430ux445}{%
\subsection{В.2 Хоёр талаас шалгах}\label{ux432.2-ux445ux43eux451ux440-ux442ux430ux43bux430ux430ux441-ux448ux430ux43bux433ux430ux445}}

\textbf{{[}ЗК · Ubuntu{]}} (цонх 3)

\begin{Shaded}
\begin{Highlighting}[]
\ExtensionTok{mosquitto\_sub}\NormalTok{ {-}h localhost {-}t }\StringTok{\textquotesingle{}cnc302/\#\textquotesingle{}}\NormalTok{ {-}v}
\end{Highlighting}
\end{Shaded}

Энэ цонх \textbf{нээлттэй хүлээж} байх ёстой --- хааж болохгүй.

\textbf{{[}Pi{]}} (цонх 1)

\begin{Shaded}
\begin{Highlighting}[]
\BuiltInTok{cd}\NormalTok{ \textasciitilde{}/cnc302/edge }\KeywordTok{\&\&} \FunctionTok{make}\NormalTok{ agent}
\end{Highlighting}
\end{Shaded}

✓ Цонх 3-т Pi-гийн телеметр 2 секунд тутам \texttt{cnc302/\allowbreak{}shutis/\allowbreak{}mhts/\allowbreak{}lab/\allowbreak{}pi3b-team\textless{}NN\textgreater{}/\allowbreak{}\ldots{}} сэдвээр гарч ирнэ. Гарвал \textbf{орчин бэлэн боллоо} --- цонх 1 ба 3-ыг \texttt{Ctrl\ +\ C}-ээр зогсоо.

\textbf{Гарч ирэхгүй бол Лаб 1 эхлэхгүй.} Дарааллаар нь шалга:

\begingroup\scriptsize\begin{longtable}[]{@{}lll@{}}
\toprule
\begin{minipage}[b]{0.30\columnwidth}\raggedright
Шалгах зүйл\strut
\end{minipage} & \begin{minipage}[b]{0.30\columnwidth}\raggedright
Команд\strut
\end{minipage} & \begin{minipage}[b]{0.30\columnwidth}\raggedright
Хүлээгдэх\strut
\end{minipage}\tabularnewline
\midrule
\endhead
\begin{minipage}[t]{0.30\columnwidth}\raggedright
1. Pi үүлийг харж байна уу\strut
\end{minipage} & \begin{minipage}[t]{0.30\columnwidth}\raggedright
{[}Pi{]} \texttt{ping\ -c3\ \textless{}LAPTOP\_\allowbreak{}IP\textgreater{}}\strut
\end{minipage} & \begin{minipage}[t]{0.30\columnwidth}\raggedright
\texttt{3\ received}\strut
\end{minipage}\tabularnewline
\begin{minipage}[t]{0.30\columnwidth}\raggedright
2. Порт нээлттэй юу\strut
\end{minipage} & \begin{minipage}[t]{0.30\columnwidth}\raggedright
{[}Pi{]} \texttt{nc\ -vz\ \textless{}LAPTOP\_\allowbreak{}IP\textgreater{}\ 1883}\strut
\end{minipage} & \begin{minipage}[t]{0.30\columnwidth}\raggedright
\texttt{succeeded} / \texttt{open}\strut
\end{minipage}\tabularnewline
\begin{minipage}[t]{0.30\columnwidth}\raggedright
3. Гүүрний тохиргоо\strut
\end{minipage} & \begin{minipage}[t]{0.30\columnwidth}\raggedright
{[}Pi{]} \texttt{docker\ compose\ logs\ mosquitto\ \textbackslash{}\textbar{}\ grep\ -i\ bridge}\strut
\end{minipage} & \begin{minipage}[t]{0.30\columnwidth}\raggedright
\texttt{Connecting\ bridge}\strut
\end{minipage}\tabularnewline
\begin{minipage}[t]{0.30\columnwidth}\raggedright
4. Сэдвийн угтвар\strut
\end{minipage} & \begin{minipage}[t]{0.30\columnwidth}\raggedright
{[}Pi{]} \texttt{grep\ \textquotesingle{}\^{}topic\textquotesingle{}\ mosquitto/conf.d/bridge.conf}\strut
\end{minipage} & \begin{minipage}[t]{0.30\columnwidth}\raggedright
\texttt{cnc302/shutis/\#}\strut
\end{minipage}\tabularnewline
\begin{minipage}[t]{0.30\columnwidth}\raggedright
5. Агентын сэдэв\strut
\end{minipage} & \begin{minipage}[t]{0.30\columnwidth}\raggedright
{[}Pi{]} \texttt{grep\ SITE\ .env}\strut
\end{minipage} & \begin{minipage}[t]{0.30\columnwidth}\raggedright
\texttt{SITE=shutis}\strut
\end{minipage}\tabularnewline
\bottomrule
\end{longtable}\endgroup

\begin{itemize}
\tightlist
\item
  1-р алхам бүтэлгүйтвэл → сүлжээ (А.5): нэг subnet мөн үү, client isolation байна уу.
\item
  2--3-р алхамд \texttt{Connection\ refused}, timeout эсвэл удаа дараа тасарч байвал → А.6 галт хана ба Private профайл; зөөврийн компьютер дээр стек асаалттай эсэх.
\item
  4 ба 5-р алхмын \texttt{SITE} зөрүүтэй бол мессеж дамжихгүй. Энэ бол \textbf{хамгийн түгээмэл алдаа}.
\end{itemize}

\hypertarget{ux432.3-vs-code-remote-ssh-ux445ux43eux43bux431ux43eux43bux442}{%
\subsection{В.3 VS Code Remote-SSH холболт}\label{ux432.3-vs-code-remote-ssh-ux445ux43eux43bux431ux43eux43bux442}}

\begin{enumerate}
\def\labelenumi{\arabic{enumi}.}
\tightlist
\item
  VS Code нээ → \texttt{F1} (эсвэл \texttt{Ctrl\ +\ Shift\ +\ P}) → \texttt{Remote-SSH:\allowbreak{}\ Connect\ to\ Host.\allowbreak{}.\allowbreak{}.\allowbreak{}} гэж бичээд сонго.
\item
  Жагсаалтаас \textbf{\texttt{pi}}-г сонго (А.7, Алхам 3-т үүсгэсэн).
\item
  Платформ асуувал \textbf{Linux}.
\item
  Анх удаа Pi дээр VS Code Server суулгана --- 3--5 мин хүлээ. Зүүн доод буланд \textbf{SSH: pi} гарна.
\item
  \textbf{File → Open Folder} → \texttt{/\allowbreak{}home/\allowbreak{}cnc302/\allowbreak{}cnc302} → \textbf{OK}.
\item
  \textbf{Terminal → New Terminal} --- энэ терминал одоо \textbf{Pi дээр} ажиллана.
\end{enumerate}

\begin{quote}
VS Code-ийн баримтаар Remote-SSH-ийн алсын хостод \textbf{1 GB RAM шаардлагатай, 2 GB RAM ба 2 цөм зөвлөмжтэй} --- Pi 3B яг доод хязгаар дээр. Дэмжигдэх платформ: Debian 64-бит ARMv8 (AArch64), kernel ≥ 4.18, glibc ≥ 2.28. Хэмжилт хийхийн өмнө \textbf{VS Code-ийн холболтыг таслах} --- эс бөгөөс санах ойн тоо гажина (\texttt{pkill\ -f\ vscode-server}).

Албан ёсны баримт: \href{https://code.visualstudio.com/docs/remote/ssh}{Remote Development using SSH --- System requirements} · \href{https://code.visualstudio.com/docs/remote/linux}{Remote Development with Linux}
\end{quote}

\hypertarget{ux433.-ux431ux44dux43bux44dux43d-ux431ux430ux439ux434ux43bux44bux43d-ux448ux430ux43bux433ux430ux43bux442}{%
\section{Г. Бэлэн байдлын шалгалт}\label{ux433.-ux431ux44dux43bux44dux43d-ux431ux430ux439ux434ux43bux44bux43d-ux448ux430ux43bux433ux430ux43bux442}}

Лаб 1-д ирэхээс өмнө эдгээр \textbf{бүгд ✔} байх ёстой. Командын гаралтыг \textbf{хуулж} Лаб 1-ийн тайлангийн хавсралтад оруулна.

\hypertarget{ux437ux43a-ux437ux4e9ux4e9ux432ux440ux438ux439ux43d-ux43aux43eux43cux43fux44cux44eux442ux435ux440}{%
\subsection{{[}ЗК{]} Зөөврийн компьютер}\label{ux437ux43a-ux437ux4e9ux4e9ux432ux440ux438ux439ux43d-ux43aux43eux43cux43fux44cux44eux442ux435ux440}}

\begin{itemize}
\tightlist
\item[$\square$]
  \textbf{{[}ЗК · PowerShell{]}} \texttt{wsl\ -\/-version} → WSL ≥ 2.1.5; \texttt{wsl\ -l\ -v} → Ubuntu, VERSION 2
\item[$\square$]
  \textbf{{[}ЗК · Ubuntu{]}} \texttt{docker\ run\ -\/\allowbreak{}-rm\ hello-world} → \texttt{Hello\ from\ Docker!}
\item[$\square$]
  \textbf{{[}ЗК · Ubuntu{]}} \texttt{docker\ info} санах ой ≈ 6 GB (А.2)
\item[$\square$]
  \textbf{{[}ЗК · Ubuntu{]}} \texttt{cd\ \textasciitilde{}/cnc302/stack\ \&\&\ make\ health} → дөрвөн мөр 200
\item[$\square$]
  \textbf{{[}ЗК · PowerShell{]}} \texttt{ipconfig} → \texttt{\textless{}LAPTOP\_\allowbreak{}IP\textgreater{}} тодорхой, сүлжээ \textbf{Private}
\item[$\square$]
  \textbf{{[}ЗК · PowerShell{]}} \texttt{Get-NetFirewallRule\ -DisplayName\ "CNC302\ cloud"} → Enabled True
\item[$\square$]
  \textbf{{[}ЗК · Ubuntu{]}} \texttt{python\ tools/\allowbreak{}sim\_\allowbreak{}device.\allowbreak{}py\ -\/\allowbreak{}-dry-run\ -\/\allowbreak{}-devices\ 2\ -\/\allowbreak{}-count\ 3} ажиллаж байна
\end{itemize}

\hypertarget{pi-raspberry-pi-3b}{%
\subsection{{[}Pi{]} Raspberry Pi 3B}\label{pi-raspberry-pi-3b}}

\begin{itemize}
\tightlist
\item[$\square$]
  \texttt{ssh\ pi} → нууц үггүй холбогдож байна
\item[$\square$]
  \texttt{uname\ -m} → \textbf{aarch64}
\item[$\square$]
  \texttt{docker\ run\ -\/\allowbreak{}-rm\ hello-world} → \texttt{sudo}-гүй ажиллаж байна
\item[$\square$]
  \texttt{free\ -h} → Swap 2.0Gi; \texttt{MemTotal}-ийг өмнө/дараа нь тэмдэглэсэн (Б.3)
\item[$\square$]
  \texttt{grep\ -w\ memory\ /\allowbreak{}sys/\allowbreak{}fs/\allowbreak{}cgroup/\allowbreak{}cgroup.\allowbreak{}controllers} → \texttt{memory} агуулна
\item[$\square$]
  \texttt{vcgencmd\ get\_\allowbreak{}throttled} → \textbf{0x0}
\item[$\square$]
  \texttt{timedatectl\ status} → synchronized: yes
\item[$\square$]
  \texttt{cd\ \textasciitilde{}/cnc302/edge\ \&\&\ make\ check} → бүгд OK
\end{itemize}

\hypertarget{ux437ux43api-ux445ux430ux43cux442ux434ux430ux430}{%
\subsection{{[}ЗК{]}{[}Pi{]} Хамтдаа}\label{ux437ux43api-ux445ux430ux43cux442ux434ux430ux430}}

\begin{itemize}
\tightlist
\item[$\square$]
  \texttt{make\ link} → \texttt{bridge/state\ 1}
\item[$\square$]
  Pi-гийн агентын телеметр зөөврийн компьютерийн \texttt{mosquitto\_sub}-д харагдаж байна (В.2)
\item[$\square$]
  Багийн Git санд эхний commit хийгдсэн, багш унших эрхтэй
\end{itemize}

\hypertarget{vm-ux43bux430ux431-8-ux430ux430ux441-ux4e9ux43cux43dux4e9}{%
\subsection{{[}VM{]} Лаб 8-аас өмнө}\label{vm-ux43bux430ux431-8-ux430ux430ux441-ux4e9ux43cux43dux4e9}}

\begin{itemize}
\tightlist
\item[$\square$]
  А.8-ын VM асаж, \texttt{ssh}-ээр холбогдож байна; \texttt{ip\ -4\ addr} → LAN-ын хаяг (10.0.2.x БИШ)
\end{itemize}

Аль нэг нь ✘ бол \textbf{лабораторийн өмнөх өдөр} багшид хандана. \textbf{Лабораторийн 4 цагийг орчин тохируулахад зарцуулж болохгүй.}

\hypertarget{ux434.-ux43eux43dux43eux448ux438ux43bux433ux43eux43e-ux442ux4afux433ux44dux44dux43cux44dux43b-ux430ux43bux434ux430ux430-ux431ux430-ux448ux438ux439ux434ux44dux43b}{%
\section{Д. Оношилгоо --- түгээмэл алдаа ба шийдэл}\label{ux434.-ux43eux43dux43eux448ux438ux43bux433ux43eux43e-ux442ux4afux433ux44dux44dux43cux44dux43b-ux430ux43bux434ux430ux430-ux431ux430-ux448ux438ux439ux434ux44dux43b}}

\hypertarget{ux437ux43a-ux437ux4e9ux4e9ux432ux440ux438ux439ux43d-ux43aux43eux43cux43fux44cux44eux442ux435ux440-1}{%
\subsection{{[}ЗК{]} Зөөврийн компьютер}\label{ux437ux43a-ux437ux4e9ux4e9ux432ux440ux438ux439ux43d-ux43aux43eux43cux43fux44cux44eux442ux435ux440-1}}

\begingroup\scriptsize\begin{longtable}[]{@{}lll@{}}
\toprule
\begin{minipage}[b]{0.30\columnwidth}\raggedright
Шинж тэмдэг\strut
\end{minipage} & \begin{minipage}[b]{0.30\columnwidth}\raggedright
Шалтгаан\strut
\end{minipage} & \begin{minipage}[b]{0.30\columnwidth}\raggedright
Шийдэл\strut
\end{minipage}\tabularnewline
\midrule
\endhead
\begin{minipage}[t]{0.30\columnwidth}\raggedright
\texttt{wsl\ -\/-install} 0.0\%-д гацах\strut
\end{minipage} & \begin{minipage}[t]{0.30\columnwidth}\raggedright
Store-оос татаж чадахгүй\strut
\end{minipage} & \begin{minipage}[t]{0.30\columnwidth}\raggedright
\texttt{wsl\ -\/\allowbreak{}-install\ -\/\allowbreak{}-web-download\ -d\ Ubuntu}\strut
\end{minipage}\tabularnewline
\begin{minipage}[t]{0.30\columnwidth}\raggedright
Docker Desktop: \texttt{WSL\ 2\ installation\ is\ incomplete} / \texttt{WSL\ update\ failed}\strut
\end{minipage} & \begin{minipage}[t]{0.30\columnwidth}\raggedright
WSL суугаагүй эсвэл хуучин\strut
\end{minipage} & \begin{minipage}[t]{0.30\columnwidth}\raggedright
А.1.3-ыг дахин хий, \texttt{wsl\ -\/-update}, restart\strut
\end{minipage}\tabularnewline
\begin{minipage}[t]{0.30\columnwidth}\raggedright
Docker Desktop: \texttt{Virtualization\ support\ not\ detected}\strut
\end{minipage} & \begin{minipage}[t]{0.30\columnwidth}\raggedright
BIOS-д VT-x/SVM унтраалттай\strut
\end{minipage} & \begin{minipage}[t]{0.30\columnwidth}\raggedright
А.1.1, Алхам 2\strut
\end{minipage}\tabularnewline
\begin{minipage}[t]{0.30\columnwidth}\raggedright
\texttt{You\ are\ not\ allowed\ to\ use\ Docker,\ you\ must\ be\ in\ the\ "docker-users"\ group}\strut
\end{minipage} & \begin{minipage}[t]{0.30\columnwidth}\raggedright
"All users" горимоор суулгасан\strut
\end{minipage} & \begin{minipage}[t]{0.30\columnwidth}\raggedright
\textbf{{[}ЗК · PowerShell (Admin){]}} \texttt{net\ localgroup\ docker-users\ \$env:\allowbreak{}USERNAME\ /\allowbreak{}add} → Windows-оос \textbf{sign out} хийж дахин нэвтэр\strut
\end{minipage}\tabularnewline
\begin{minipage}[t]{0.30\columnwidth}\raggedright
\texttt{docker:\allowbreak{}\ error\ during\ connect} / \texttt{Cannot\ connect\ to\ the\ Docker\ daemon}\strut
\end{minipage} & \begin{minipage}[t]{0.30\columnwidth}\raggedright
Docker Desktop ажиллаагүй\strut
\end{minipage} & \begin{minipage}[t]{0.30\columnwidth}\raggedright
Docker Desktop нээж, \textbf{Engine running} болтол хүлээ\strut
\end{minipage}\tabularnewline
\begin{minipage}[t]{0.30\columnwidth}\raggedright
Ubuntu-д \texttt{docker:\allowbreak{}\ command\ not\ found}\strut
\end{minipage} & \begin{minipage}[t]{0.30\columnwidth}\raggedright
WSL integration асаагүй\strut
\end{minipage} & \begin{minipage}[t]{0.30\columnwidth}\raggedright
А.1.6, Алхам 3\strut
\end{minipage}\tabularnewline
\begin{minipage}[t]{0.30\columnwidth}\raggedright
Docker Desktop гацах, компьютер удаашрах\strut
\end{minipage} & \begin{minipage}[t]{0.30\columnwidth}\raggedright
RAM хүрэлцэхгүй\strut
\end{minipage} & \begin{minipage}[t]{0.30\columnwidth}\raggedright
А.2-т \texttt{memory}-г 1 GB-аар багасга; Chrome-ын табуудыг хаа\strut
\end{minipage}\tabularnewline
\begin{minipage}[t]{0.30\columnwidth}\raggedright
\texttt{docker\ info} санах ой 50\% хэвээр\strut
\end{minipage} & \begin{minipage}[t]{0.30\columnwidth}\raggedright
\texttt{.wslconfig} буруу нэртэй/газарт\strut
\end{minipage} & \begin{minipage}[t]{0.30\columnwidth}\raggedright
А.2-ын ⚠; \texttt{wsl\ -\/\allowbreak{}-shutdown} хийсэн эсэх\strut
\end{minipage}\tabularnewline
\begin{minipage}[t]{0.30\columnwidth}\raggedright
\texttt{Bind\ for\ 0.\allowbreak{}0.\allowbreak{}0.\allowbreak{}0:\allowbreak{}1883\ failed:\allowbreak{}\ port\ is\ already\ allocated}\strut
\end{minipage} & \begin{minipage}[t]{0.30\columnwidth}\raggedright
Windows дээр өөр MQTT брокер (Mosquitto service) ажиллаж байна\strut
\end{minipage} & \begin{minipage}[t]{0.30\columnwidth}\raggedright
\texttt{Win\ +\ R} → \texttt{services.msc} → \textbf{Mosquitto Broker} → \textbf{Stop}, Startup type: \textbf{Disabled}\strut
\end{minipage}\tabularnewline
\begin{minipage}[t]{0.30\columnwidth}\raggedright
\texttt{pull\ access\ denied\ for\ cnc302/\allowbreak{}registry}\strut
\end{minipage} & \begin{minipage}[t]{0.30\columnwidth}\raggedright
\texttt{-\/\allowbreak{}-ignore-buildable}-гүй \texttt{pull}\strut
\end{minipage} & \begin{minipage}[t]{0.30\columnwidth}\raggedright
Анхааралгүй байж болно --- \texttt{make\ up} өөрөө барина. Эсвэл А.4, Алхам 2\strut
\end{minipage}\tabularnewline
\begin{minipage}[t]{0.30\columnwidth}\raggedright
\texttt{pull} маш удаан, тасраад байх\strut
\end{minipage} & \begin{minipage}[t]{0.30\columnwidth}\raggedright
Интернэт удаан\strut
\end{minipage} & \begin{minipage}[t]{0.30\columnwidth}\raggedright
Дахин ажиллуул --- татсан давхаргуудаа үргэлжлүүлнэ\strut
\end{minipage}\tabularnewline
\begin{minipage}[t]{0.30\columnwidth}\raggedright
\texttt{make\ health}-д \texttt{000} эсвэл \texttt{DOWN}\strut
\end{minipage} & \begin{minipage}[t]{0.30\columnwidth}\raggedright
Контейнер асаагүй / асаж байна\strut
\end{minipage} & \begin{minipage}[t]{0.30\columnwidth}\raggedright
1 мин хүлээ; \texttt{make\ status}; \texttt{make\ logs\ S=\allowbreak{}emqx}\strut
\end{minipage}\tabularnewline
\begin{minipage}[t]{0.30\columnwidth}\raggedright
EMQX/Grafana нууц үг таарахгүй\strut
\end{minipage} & \begin{minipage}[t]{0.30\columnwidth}\raggedright
\texttt{.env}-ийг стек анх асаасны дараа сольсон\strut
\end{minipage} & \begin{minipage}[t]{0.30\columnwidth}\raggedright
Өгөгдлийг устгаж дахин эхлэх: \texttt{make\ clean} (\texttt{yes}) → \texttt{make\ up}\strut
\end{minipage}\tabularnewline
\begin{minipage}[t]{0.30\columnwidth}\raggedright
\texttt{/bin/sh\^{}M:\ bad\ interpreter}\strut
\end{minipage} & \begin{minipage}[t]{0.30\columnwidth}\raggedright
\texttt{.sh} файлд CRLF орсон\strut
\end{minipage} & \begin{minipage}[t]{0.30\columnwidth}\raggedright
А.3, Алхам 1-ийн \texttt{core.\allowbreak{}autocrlf\ input} → сангаа устгаж \textbf{дахин clone} хий\strut
\end{minipage}\tabularnewline
\begin{minipage}[t]{0.30\columnwidth}\raggedright
\texttt{git\ clone} нэвтрэлт асуугаад бүтэлгүйтэх\strut
\end{minipage} & \begin{minipage}[t]{0.30\columnwidth}\raggedright
Credential Manager тохируулаагүй\strut
\end{minipage} & \begin{minipage}[t]{0.30\columnwidth}\raggedright
А.3, Алхам 1-ийн \texttt{credential.\allowbreak{}helper} мөр\strut
\end{minipage}\tabularnewline
\bottomrule
\end{longtable}\endgroup

\hypertarget{pi-raspberry-pi}{%
\subsection{{[}Pi{]} Raspberry Pi}\label{pi-raspberry-pi}}

\begingroup\scriptsize\begin{longtable}[]{@{}lll@{}}
\toprule
\begin{minipage}[b]{0.30\columnwidth}\raggedright
Шинж тэмдэг\strut
\end{minipage} & \begin{minipage}[b]{0.30\columnwidth}\raggedright
Шалтгаан\strut
\end{minipage} & \begin{minipage}[b]{0.30\columnwidth}\raggedright
Шийдэл\strut
\end{minipage}\tabularnewline
\midrule
\endhead
\begin{minipage}[t]{0.30\columnwidth}\raggedright
\texttt{Ping\ request\ could\ not\ find\ host\ pi-teamNN.\allowbreak{}local}\strut
\end{minipage} & \begin{minipage}[t]{0.30\columnwidth}\raggedright
Pi асаж дуусаагүй / mDNS ажиллахгүй / Pi сүлжээнд ороогүй\strut
\end{minipage} & \begin{minipage}[t]{0.30\columnwidth}\raggedright
2 мин хүлээ; router-ийн жагсаалтаас IP-г ол; \texttt{\textless{}PI\_\allowbreak{}IP\textgreater{}}-ээр холбогд\strut
\end{minipage}\tabularnewline
\begin{minipage}[t]{0.30\columnwidth}\raggedright
Pi сүлжээнд огт гарч ирэхгүй\strut
\end{minipage} & \begin{minipage}[t]{0.30\columnwidth}\raggedright
Imager-ийн customisation хэрэгжээгүй, Wi-Fi SSID/нууц үг буруу, 5 GHz сүлжээ\strut
\end{minipage} & \begin{minipage}[t]{0.30\columnwidth}\raggedright
Кабелиар холбо; Imager-ийг дахин нээж картаа \textbf{дахин бич}; HDMI дэлгэц + гар залгаж шалгаж болно\strut
\end{minipage}\tabularnewline
\begin{minipage}[t]{0.30\columnwidth}\raggedright
\texttt{Permission\ denied\ (publickey)}\strut
\end{minipage} & \begin{minipage}[t]{0.30\columnwidth}\raggedright
Imager-т өөр түлхүүр буулгасан, эсвэл \texttt{ssh} буруу түлхүүр ашиглаж байна\strut
\end{minipage} & \begin{minipage}[t]{0.30\columnwidth}\raggedright
\texttt{ssh\ -i\ "\$env:USERPROFILE\textbackslash{}.ssh\textbackslash{}cnc302"\ cnc302@\textless{}PI\_IP\textgreater{}}; хэрэглэгчийн нэр \texttt{cnc302} эсэх (\texttt{pi@} биш); болохгүй бол картаа А.7-ийн \texttt{.pub}-аар \textbf{дахин бич}\strut
\end{minipage}\tabularnewline
\begin{minipage}[t]{0.30\columnwidth}\raggedright
\texttt{WARNING:\allowbreak{}\ REMOTE\ HOST\ IDENTIFICATION\ HAS\ CHANGED!}\strut
\end{minipage} & \begin{minipage}[t]{0.30\columnwidth}\raggedright
Картаа дахин бичсэн тул Pi-гийн түлхүүр солигдсон\strut
\end{minipage} & \begin{minipage}[t]{0.30\columnwidth}\raggedright
\textbf{{[}ЗК · PowerShell{]}} \texttt{ssh-keygen\ -R\ pi-team\textless{}NN\textgreater{}.\allowbreak{}local} (ба \texttt{ssh-keygen\ -R\ \textless{}PI\_\allowbreak{}IP\textgreater{}}) → дахин холбогд\strut
\end{minipage}\tabularnewline
\begin{minipage}[t]{0.30\columnwidth}\raggedright
\texttt{uname\ -m} → \texttt{armv7l}\strut
\end{minipage} & \begin{minipage}[t]{0.30\columnwidth}\raggedright
32-bit OS бичсэн\strut
\end{minipage} & \begin{minipage}[t]{0.30\columnwidth}\raggedright
Б.1 --- \textbf{Lite (64-bit)}\strut
\end{minipage}\tabularnewline
\begin{minipage}[t]{0.30\columnwidth}\raggedright
\texttt{Suites:} хоосон / \texttt{apt\ update}-д Docker-ийн 404\strut
\end{minipage} & \begin{minipage}[t]{0.30\columnwidth}\raggedright
Б.4, Алхам 3-ыг хэсэгчлэн хуулсан\strut
\end{minipage} & \begin{minipage}[t]{0.30\columnwidth}\raggedright
Алхам 3-ыг блокоор нь дахин буулга\strut
\end{minipage}\tabularnewline
\begin{minipage}[t]{0.30\columnwidth}\raggedright
\texttt{permission\ denied\ .\allowbreak{}.\allowbreak{}.\allowbreak{}\ docker.\allowbreak{}sock}\strut
\end{minipage} & \begin{minipage}[t]{0.30\columnwidth}\raggedright
\texttt{docker} бүлэгт нэмсний дараа дахин нэвтрээгүй\strut
\end{minipage} & \begin{minipage}[t]{0.30\columnwidth}\raggedright
\texttt{exit} → \texttt{ssh\ pi}\strut
\end{minipage}\tabularnewline
\begin{minipage}[t]{0.30\columnwidth}\raggedright
\texttt{exec\ format\ error}\strut
\end{minipage} & \begin{minipage}[t]{0.30\columnwidth}\raggedright
arm64 биш (x86) image эсвэл 32-bit OS\strut
\end{minipage} & \begin{minipage}[t]{0.30\columnwidth}\raggedright
\texttt{uname\ -m} → \texttt{aarch64} байх ёстой\strut
\end{minipage}\tabularnewline
\begin{minipage}[t]{0.30\columnwidth}\raggedright
\texttt{get\_throttled} ≠ \texttt{0x0}, Pi санамсаргүй restart хийх\strut
\end{minipage} & \begin{minipage}[t]{0.30\columnwidth}\raggedright
Тэжээл сул (under-voltage)\strut
\end{minipage} & \begin{minipage}[t]{0.30\columnwidth}\raggedright
Албан ёсны 5 V / 2.5 A адаптер, богино зузаан кабель\strut
\end{minipage}\tabularnewline
\begin{minipage}[t]{0.30\columnwidth}\raggedright
\texttt{free\ -h} → Swap 2.0Gi биш\strut
\end{minipage} & \begin{minipage}[t]{0.30\columnwidth}\raggedright
Reboot хийгээгүй / буруу зам\strut
\end{minipage} & \begin{minipage}[t]{0.30\columnwidth}\raggedright
Б.5, Алхам 1-ээс дахин; \texttt{sudo\ reboot}\strut
\end{minipage}\tabularnewline
\begin{minipage}[t]{0.30\columnwidth}\raggedright
\texttt{make\ check} → \texttt{cgroup\ mem\ ✗}\strut
\end{minipage} & \begin{minipage}[t]{0.30\columnwidth}\raggedright
\texttt{cmdline.txt} хоёр мөр болсон эсвэл reboot хийгээгүй\strut
\end{minipage} & \begin{minipage}[t]{0.30\columnwidth}\raggedright
\texttt{cat\ /\allowbreak{}boot/\allowbreak{}firmware/\allowbreak{}cmdline.\allowbreak{}txt} нэг мөр эсэх; reboot\strut
\end{minipage}\tabularnewline
\begin{minipage}[t]{0.30\columnwidth}\raggedright
\texttt{pip\ install} маш удаан, "гацсан мэт"\strut
\end{minipage} & \begin{minipage}[t]{0.30\columnwidth}\raggedright
Pi 3B удаан\strut
\end{minipage} & \begin{minipage}[t]{0.30\columnwidth}\raggedright
15 мин хүлээ; \texttt{htop}-оор CPU ачаалал байгаа эсэхийг шалга\strut
\end{minipage}\tabularnewline
\begin{minipage}[t]{0.30\columnwidth}\raggedright
\texttt{bridge/state\ 0}, Pi → 1883 холбогдохгүй\strut
\end{minipage} & \begin{minipage}[t]{0.30\columnwidth}\raggedright
(а) стек асаагүй (б) Windows Firewall (в) Public профайл (г) client isolation (д) \texttt{CLOUD\_HOST} буруу\strut
\end{minipage} & \begin{minipage}[t]{0.30\columnwidth}\raggedright
(а) \texttt{make\ health} (б) А.6 (в) А.5 (г) лабын router/hotspot (д) \texttt{grep\ CLOUD\_HOST\ \textasciitilde{}/cnc302/edge/.env} → \texttt{make\ bridge} → \texttt{make\ restart}\strut
\end{minipage}\tabularnewline
\bottomrule
\end{longtable}\endgroup

Илүү олон тохиолдол: \texttt{docs/\allowbreak{}troubleshooting.\allowbreak{}md}.

\textbf{Pi-г унтраахдаа} тэжээлийг шууд бүү сугал --- microSD эвдэрнэ. \texttt{sudo\ poweroff} гэж бичээд \textbf{ногоон гэрэл анивчихаа больсны дараа} сугал.

\hypertarget{ux44dux445-ux441ux443ux440ux432ux430ux43bux436}{%
\section{Эх сурвалж}\label{ux44dux445-ux441ux443ux440ux432ux430ux43bux436}}

\begingroup\scriptsize\begin{longtable}[]{@{}llll@{}}
\toprule
\begin{minipage}[b]{0.22\columnwidth}\raggedright
\#\strut
\end{minipage} & \begin{minipage}[b]{0.22\columnwidth}\raggedright
Эх сурвалж\strut
\end{minipage} & \begin{minipage}[b]{0.22\columnwidth}\raggedright
Юуг баталгаажуулсан\strut
\end{minipage} & \begin{minipage}[b]{0.22\columnwidth}\raggedright
Хандсан\strut
\end{minipage}\tabularnewline
\midrule
\endhead
\begin{minipage}[t]{0.22\columnwidth}\raggedright
1\strut
\end{minipage} & \begin{minipage}[t]{0.22\columnwidth}\raggedright
\href{https://docs.docker.com/desktop/setup/install/windows-install/}{Install Docker Desktop on Windows}\strut
\end{minipage} & \begin{minipage}[t]{0.22\columnwidth}\raggedright
WSL 2 backend: Windows 10 22H2 (19045) / Windows 11 23H2 (22631)+, WSL ≥ 2.1.5, 8 GB RAM, virtualization; Home нь зөвхөн Linux контейнер; "Use WSL 2 instead of Hyper-V"; Subscription Service Agreement; \texttt{docker-users} бүлэг\strut
\end{minipage} & \begin{minipage}[t]{0.22\columnwidth}\raggedright
2026-10\strut
\end{minipage}\tabularnewline
\begin{minipage}[t]{0.22\columnwidth}\raggedright
2\strut
\end{minipage} & \begin{minipage}[t]{0.22\columnwidth}\raggedright
\href{https://docs.docker.com/desktop/features/wsl/}{Docker Desktop --- WSL 2 backend}, \href{https://docs.docker.com/desktop/features/wsl/best-practices/}{Best practices}\strut
\end{minipage} & \begin{minipage}[t]{0.22\columnwidth}\raggedright
Use WSL 2 based engine; Resources → WSL Integration; \texttt{wsl\ -l\ -v}, \texttt{-\/-set-version}; кодыг Linux файлын системд хадгалах\strut
\end{minipage} & \begin{minipage}[t]{0.22\columnwidth}\raggedright
2026-10\strut
\end{minipage}\tabularnewline
\begin{minipage}[t]{0.22\columnwidth}\raggedright
3\strut
\end{minipage} & \begin{minipage}[t]{0.22\columnwidth}\raggedright
\href{https://learn.microsoft.com/en-us/windows/wsl/install}{Microsoft --- Install WSL}\strut
\end{minipage} & \begin{minipage}[t]{0.22\columnwidth}\raggedright
\texttt{wsl\ -\/\allowbreak{}-install\ -d}, \texttt{-\/-web-download}, анх асаахад UNIX хэрэглэгч үүсгэх, \texttt{wsl\ -l\ -v}\strut
\end{minipage} & \begin{minipage}[t]{0.22\columnwidth}\raggedright
2026-10\strut
\end{minipage}\tabularnewline
\begin{minipage}[t]{0.22\columnwidth}\raggedright
4\strut
\end{minipage} & \begin{minipage}[t]{0.22\columnwidth}\raggedright
\href{https://github.com/git-ecosystem/git-credential-manager/blob/main/docs/wsl.md}{Git Credential Manager --- WSL}\strut
\end{minipage} & \begin{minipage}[t]{0.22\columnwidth}\raggedright
WSL-ийн Git-ийг Windows-ийн GCM-тэй холбох \texttt{credential.\allowbreak{}helper}\strut
\end{minipage} & \begin{minipage}[t]{0.22\columnwidth}\raggedright
2026-10\strut
\end{minipage}\tabularnewline
\begin{minipage}[t]{0.22\columnwidth}\raggedright
5\strut
\end{minipage} & \begin{minipage}[t]{0.22\columnwidth}\raggedright
\href{https://docs.docker.com/desktop/settings-and-maintenance/settings/}{Docker Desktop --- Settings}\strut
\end{minipage} & \begin{minipage}[t]{0.22\columnwidth}\raggedright
Memory limit: Mac, Linux, Windows Hyper-V; анхдагч нь хостын 50\%; WSL 2 горимд санах ойг WSL 2 VM дээр тохируулна\strut
\end{minipage} & \begin{minipage}[t]{0.22\columnwidth}\raggedright
2026-09\strut
\end{minipage}\tabularnewline
\begin{minipage}[t]{0.22\columnwidth}\raggedright
6\strut
\end{minipage} & \begin{minipage}[t]{0.22\columnwidth}\raggedright
\href{https://learn.microsoft.com/en-us/windows/wsl/wsl-config}{WSL --- Advanced settings configuration (.\allowbreak{}wslconfig)}\strut
\end{minipage} & \begin{minipage}[t]{0.22\columnwidth}\raggedright
\texttt{\%UserProfile\%\textbackslash{}.wslconfig}, \texttt{{[}wsl2{]}\ memory=\allowbreak{}}, анхдагч 50\%, \texttt{wsl\ -\/\allowbreak{}-shutdown}\strut
\end{minipage} & \begin{minipage}[t]{0.22\columnwidth}\raggedright
2026-09\strut
\end{minipage}\tabularnewline
\begin{minipage}[t]{0.22\columnwidth}\raggedright
7\strut
\end{minipage} & \begin{minipage}[t]{0.22\columnwidth}\raggedright
\href{https://learn.microsoft.com/en-us/powershell/module/netsecurity/new-netfirewallrule}{New-NetFirewallRule}\strut
\end{minipage} & \begin{minipage}[t]{0.22\columnwidth}\raggedright
\texttt{-LocalPort} порт/жагсаалт, \texttt{-Profile\ Private}\strut
\end{minipage} & \begin{minipage}[t]{0.22\columnwidth}\raggedright
2026-09\strut
\end{minipage}\tabularnewline
\begin{minipage}[t]{0.22\columnwidth}\raggedright
8\strut
\end{minipage} & \begin{minipage}[t]{0.22\columnwidth}\raggedright
\href{https://www.virtualbox.org/manual/UserManual.html}{VirtualBox User Manual}\strut
\end{minipage} & \begin{minipage}[t]{0.22\columnwidth}\raggedright
Hyper-V ажиллаж буй хост дээр гүйцэтгэл буурна, Windows Hypervisor Platform шаардлагатай; Bridged networking\strut
\end{minipage} & \begin{minipage}[t]{0.22\columnwidth}\raggedright
2026-09\strut
\end{minipage}\tabularnewline
\begin{minipage}[t]{0.22\columnwidth}\raggedright
9\strut
\end{minipage} & \begin{minipage}[t]{0.22\columnwidth}\raggedright
\href{https://ubuntu.com/about/release-cycle}{Ubuntu release cycle}\strut
\end{minipage} & \begin{minipage}[t]{0.22\columnwidth}\raggedright
Одоогийн LTS: 26.04 (24.04 LTS дэмжигдсэн)\strut
\end{minipage} & \begin{minipage}[t]{0.22\columnwidth}\raggedright
2026-09\strut
\end{minipage}\tabularnewline
\begin{minipage}[t]{0.22\columnwidth}\raggedright
10\strut
\end{minipage} & \begin{minipage}[t]{0.22\columnwidth}\raggedright
\href{https://docs.k3s.io/installation/requirements}{K3s --- Requirements}\strut
\end{minipage} & \begin{minipage}[t]{0.22\columnwidth}\raggedright
server 2 цөм/2 GB, agent 1 цөм/512 MB; Raspberry Pi OS-д \texttt{cgroup\_\allowbreak{}memory=\allowbreak{}1\ cgroup\_\allowbreak{}enable=\allowbreak{}memory}, \texttt{/\allowbreak{}boot/\allowbreak{}firmware/\allowbreak{}cmdline.\allowbreak{}txt}\strut
\end{minipage} & \begin{minipage}[t]{0.22\columnwidth}\raggedright
2026-09\strut
\end{minipage}\tabularnewline
\begin{minipage}[t]{0.22\columnwidth}\raggedright
11\strut
\end{minipage} & \begin{minipage}[t]{0.22\columnwidth}\raggedright
\href{https://www.raspberrypi.com/documentation/computers/raspberry-pi.html}{Raspberry Pi hardware}, \href{https://www.raspberrypi.com/documentation/computers/processors.html\#bcm2837}{BCM2837}\strut
\end{minipage} & \begin{minipage}[t]{0.22\columnwidth}\raggedright
Pi 3B: BCM2837, 4×Cortex-A53 @1.2 GHz, 1 GB, 100 Mb/s Ethernet, 4×USB 2.0, 2.4 GHz Wi-Fi; 80--85 °C-д throttle\strut
\end{minipage} & \begin{minipage}[t]{0.22\columnwidth}\raggedright
2026-09\strut
\end{minipage}\tabularnewline
\begin{minipage}[t]{0.22\columnwidth}\raggedright
12\strut
\end{minipage} & \begin{minipage}[t]{0.22\columnwidth}\raggedright
\href{https://www.raspberrypi.com/documentation/computers/getting-started.html}{Getting started}, \href{https://www.raspberrypi.com/news/a-new-raspberry-pi-imager/}{A new Raspberry Pi Imager}\strut
\end{minipage} & \begin{minipage}[t]{0.22\columnwidth}\raggedright
Pi 3: 5 V / 2.5 A, 12.5 W; Imager 2.0-ийн алхмууд: device → OS → storage → customisation (hostname, localisation, user, Wi-Fi, remote access/SSH public key, Raspberry Pi Connect) → write\strut
\end{minipage} & \begin{minipage}[t]{0.22\columnwidth}\raggedright
2026-10\strut
\end{minipage}\tabularnewline
\begin{minipage}[t]{0.22\columnwidth}\raggedright
13\strut
\end{minipage} & \begin{minipage}[t]{0.22\columnwidth}\raggedright
\href{https://www.raspberrypi.com/documentation/computers/os.html\#vcgencmd}{vcgencmd get\_\allowbreak{}throttled}\strut
\end{minipage} & \begin{minipage}[t]{0.22\columnwidth}\raggedright
бит 0 = undervoltage, бит 2 = throttled, бит 16/18 = өмнө нь болсон\strut
\end{minipage} & \begin{minipage}[t]{0.22\columnwidth}\raggedright
2026-09\strut
\end{minipage}\tabularnewline
\begin{minipage}[t]{0.22\columnwidth}\raggedright
14\strut
\end{minipage} & \begin{minipage}[t]{0.22\columnwidth}\raggedright
\href{https://www.raspberrypi.com/documentation/computers/os.html}{Raspberry Pi OS}\strut
\end{minipage} & \begin{minipage}[t]{0.22\columnwidth}\raggedright
Сүүлийн хувилбар Trixie (өмнөх Bookworm); Lite = command-line-only; 64-бит нь Pi 3-д\strut
\end{minipage} & \begin{minipage}[t]{0.22\columnwidth}\raggedright
2026-09\strut
\end{minipage}\tabularnewline
\begin{minipage}[t]{0.22\columnwidth}\raggedright
15\strut
\end{minipage} & \begin{minipage}[t]{0.22\columnwidth}\raggedright
\href{https://www.raspberrypi.com/documentation/computers/legacy_config_txt.html}{Legacy config.\allowbreak{}txt --- gpu\_\allowbreak{}mem}\strut
\end{minipage} & \begin{minipage}[t]{0.22\columnwidth}\raggedright
1 GB-д анхдагч 76; legacy --- Bookworm+ дээр ажиллахгүй, албан ёсоор дэмжигдэхгүй; доод утга 16\strut
\end{minipage} & \begin{minipage}[t]{0.22\columnwidth}\raggedright
2026-09\strut
\end{minipage}\tabularnewline
\begin{minipage}[t]{0.22\columnwidth}\raggedright
16\strut
\end{minipage} & \begin{minipage}[t]{0.22\columnwidth}\raggedright
\href{https://www.raspberrypi.com/documentation/computers/configuration.html}{Kernel command line}\strut
\end{minipage} & \begin{minipage}[t]{0.22\columnwidth}\raggedright
\texttt{cmdline.txt}-ийн бүх параметр нэг мөрөнд\strut
\end{minipage} & \begin{minipage}[t]{0.22\columnwidth}\raggedright
2026-09\strut
\end{minipage}\tabularnewline
\begin{minipage}[t]{0.22\columnwidth}\raggedright
17\strut
\end{minipage} & \begin{minipage}[t]{0.22\columnwidth}\raggedright
\href{https://github.com/raspberrypi/rpi-swap}{raspberrypi/rpi-swap}\strut
\end{minipage} & \begin{minipage}[t]{0.22\columnwidth}\raggedright
\texttt{dphys-swapfile}-ийг орлоно; \texttt{/\allowbreak{}etc/\allowbreak{}rpi/\allowbreak{}swap.\allowbreak{}conf.\allowbreak{}d/\allowbreak{}}; \texttt{Mechanism=auto} = zram+file; Example 1 (\texttt{swapfile}, \texttt{FixedSizeMiB=\allowbreak{}2048}); өөрчлөлтийн дараа reboot\strut
\end{minipage} & \begin{minipage}[t]{0.22\columnwidth}\raggedright
2026-09\strut
\end{minipage}\tabularnewline
\begin{minipage}[t]{0.22\columnwidth}\raggedright
18\strut
\end{minipage} & \begin{minipage}[t]{0.22\columnwidth}\raggedright
\href{https://docs.docker.com/engine/install/debian/}{Install Docker Engine on Debian}\strut
\end{minipage} & \begin{minipage}[t]{0.22\columnwidth}\raggedright
Trixie 13 / Bookworm 12, arm64; apt repository-ийн алхам; convenience script зөвхөн туршилт/хөгжүүлэлтэд\strut
\end{minipage} & \begin{minipage}[t]{0.22\columnwidth}\raggedright
2026-09\strut
\end{minipage}\tabularnewline
\begin{minipage}[t]{0.22\columnwidth}\raggedright
19\strut
\end{minipage} & \begin{minipage}[t]{0.22\columnwidth}\raggedright
\href{https://docs.docker.com/engine/install/raspberry-pi-os/}{Install Docker Engine on Raspberry Pi OS (32-bit)}\strut
\end{minipage} & \begin{minipage}[t]{0.22\columnwidth}\raggedright
Engine v28 нь armhf-ийн сүүлийнх; 64-бит ARM-д Debian arm64 багц\strut
\end{minipage} & \begin{minipage}[t]{0.22\columnwidth}\raggedright
2026-09\strut
\end{minipage}\tabularnewline
\begin{minipage}[t]{0.22\columnwidth}\raggedright
20\strut
\end{minipage} & \begin{minipage}[t]{0.22\columnwidth}\raggedright
\href{https://docs.docker.com/engine/install/linux-postinstall/}{Linux post-installation steps}\strut
\end{minipage} & \begin{minipage}[t]{0.22\columnwidth}\raggedright
\texttt{usermod\ -aG\ docker}, дахин нэвтрэх, docker бүлгийн эрсдэл\strut
\end{minipage} & \begin{minipage}[t]{0.22\columnwidth}\raggedright
2026-09\strut
\end{minipage}\tabularnewline
\begin{minipage}[t]{0.22\columnwidth}\raggedright
21\strut
\end{minipage} & \begin{minipage}[t]{0.22\columnwidth}\raggedright
\href{https://docs.docker.com/engine/logging/drivers/json-file/}{JSON File logging driver}\strut
\end{minipage} & \begin{minipage}[t]{0.22\columnwidth}\raggedright
\texttt{daemon.json}-ийн \texttt{log-opts} мөр утгатай; зөвхөн шинэ контейнерт\strut
\end{minipage} & \begin{minipage}[t]{0.22\columnwidth}\raggedright
2026-09\strut
\end{minipage}\tabularnewline
\begin{minipage}[t]{0.22\columnwidth}\raggedright
22\strut
\end{minipage} & \begin{minipage}[t]{0.22\columnwidth}\raggedright
\href{https://code.visualstudio.com/docs/remote/ssh}{VS Code --- Remote Development using SSH}, \href{https://code.visualstudio.com/docs/remote/linux}{with Linux}\strut
\end{minipage} & \begin{minipage}[t]{0.22\columnwidth}\raggedright
1 GB RAM шаардлагатай, 2 GB + 2 цөм зөвлөмжтэй; Debian 64-бит ARMv8, kernel ≥ 4.18, glibc ≥ 2.28\strut
\end{minipage} & \begin{minipage}[t]{0.22\columnwidth}\raggedright
2026-09\strut
\end{minipage}\tabularnewline
\begin{minipage}[t]{0.22\columnwidth}\raggedright
23\strut
\end{minipage} & \begin{minipage}[t]{0.22\columnwidth}\raggedright
\href{https://pypi.org/project/ai-edge-litert/}{pypi: ai-edge-litert}, \href{https://pypi.org/project/numpy/}{pypi: numpy}\strut
\end{minipage} & \begin{minipage}[t]{0.22\columnwidth}\raggedright
ai-edge-litert: cp311/cp313 \texttt{manylinux\_\allowbreak{}2\_\allowbreak{}27\_\allowbreak{}aarch64}; numpy aarch64 wheel\strut
\end{minipage} & \begin{minipage}[t]{0.22\columnwidth}\raggedright
2026-09\strut
\end{minipage}\tabularnewline
\bottomrule
\end{longtable}\endgroup
